\documentclass[11pt]{article}
\usepackage[margin=1in]{geometry}
\usepackage{amsmath,amssymb,amsthm,mathtools}
\usepackage[nopatch=footnote]{microtype}
\usepackage{setspace}
\usepackage{enumitem}
\usepackage{booktabs}
\usepackage{longtable}
\usepackage{array}
\usepackage{multicol}
\usepackage{hyperref}
\usepackage{xcolor}
\usepackage{physics}
\usepackage{mathrsfs}
\usepackage{fancyhdr}
\usepackage{titlesec}
\usepackage{tocloft}

\hypersetup{colorlinks=true,linkcolor=blue!45!black,citecolor=blue!45!black,urlcolor=blue!45!black}
\setstretch{1.06}
\setlength{\parindent}{1.3em}
\setlength{\parskip}{0pt}
\setlength{\headheight}{14pt}
\setlist[itemize]{topsep=0.25em,itemsep=0.15em}
\setlist[enumerate]{topsep=0.25em,itemsep=0.15em}
\setcounter{tocdepth}{2}
\numberwithin{equation}{section}
\titleformat{\section}{\large\normalfont}{\thesection.}{0.6em}{}
\titleformat{\subsection}{\normalsize\normalfont\itshape}{\thesubsection.}{0.6em}{}
\titleformat{\subsubsection}[runin]{\normalfont\itshape}{}{0pt}{}[.]
\titlespacing*{\subsubsection}{0pt}{1.2ex plus .3ex}{0.6em}
\setcounter{secnumdepth}{2}
\renewcommand{\cfttoctitlefont}{\large\normalfont}
\renewcommand{\cftsecfont}{\normalfont}
\renewcommand{\cftsecpagefont}{\normalfont}
\setlength{\cftbeforesecskip}{4pt}
\renewcommand{\abstractname}{\normalfont Abstract}
\emergencystretch=1em
\widowpenalty=10000
\clubpenalty=10000
\pagestyle{fancy}
\fancyhf{}
\fancyhead[L]{SYK for Mathematicians}
\fancyhead[R]{\thepage}
\renewcommand{\headrulewidth}{0.3pt}

\newtheoremstyle{result}{6pt}{6pt}{\itshape}{}{\normalfont}{.}{.5em}{}
\newtheoremstyle{exposition}{6pt}{6pt}{\normalfont}{}{\itshape}{.}{.5em}{}
\theoremstyle{exposition}
\newtheorem{definition}{Definition}[section]
\newtheorem{remark}[definition]{Remark}
\newtheorem{example}[definition]{Example}
\newtheorem{heuristic}[definition]{Formal statement}
\theoremstyle{result}
\newtheorem{proposition}[definition]{Proposition}
\newtheorem{theorem}[definition]{Theorem}
\newtheorem{lemma}[definition]{Lemma}
\newtheorem{conjecture}[definition]{Conjecture}

\newcommand{\E}{\mathbb E}
\newcommand{\Pp}{\mathbb P}
\newcommand{\C}{\mathbb C}
\newcommand{\R}{\mathbb R}
\newcommand{\Pf}{\operatorname{Pf}}
\newcommand{\sgn}{\operatorname{sgn}}
\newcommand{\Diff}{\operatorname{Diff}}
\newcommand{\AdS}{\mathrm{AdS}}
\newcommand{\SYK}{\mathrm{SYK}}
\newcommand{\IR}{\mathrm{IR}}
\newcommand{\UV}{\mathrm{UV}}
\newcommand{\OTOC}{\mathrm{OTOC}}
\newcommand{\Sch}{\mathrm{Sch}}
\newcommand{\cH}{\mathcal H}
\newcommand{\cF}{\mathcal F}
\newcommand{\cO}{\mathcal O}
\newcommand{\cA}{\mathcal A}
\newcommand{\cG}{\mathcal G}

\title{The Sachdev--Ye--Kitaev model\\[0.3em]
\large Lecture notes for mathematicians}
\author{}
\date{September 2026}

\begin{document}
\maketitle

\begin{abstract}
The Sachdev--Ye--Kitaev model is a Gaussian random Hamiltonian generated by a Clifford algebra. These notes introduce its finite-dimensional formulation and the large-$N$ equations for its thermal correlation functions. We develop the fermionic algebra, explain the bilocal description, and derive the conformal scaling predicted by the infrared equations. We then discuss the Schwarzian effective action, four-point functions, quantum chaos, and the relation to nearly-$\AdS_2$ gravity. The intended reader knows probability, linear algebra, and elementary quantum mechanics. Finite-$N$ results are proved where elementary; diagrammatic and effective-theory arguments are accompanied by their assumptions. The main physics reference is Rosenhaus~\cite{Rosenhaus}.
\end{abstract}

\begin{multicols}{2}
\small
\tableofcontents
\end{multicols}
\newpage

\section[The finite-dimensional model]{The model and its finite-dimensional formulation}\label{sec:foundations}

\subsection{Introduction}

The Sachdev-Ye-Kitaev model is a random quantum many-body system with no spatial coordinate. In its most common form it consists of $N$ Majorana fermions and random all-to-all quartic interactions. For even interaction order $q$, the Hamiltonian is
\begin{equation}
H_q=\sigma_N\sum_{I\in\mathcal I_N}g_I\psi_I,
\qquad
\{\psi_i,\psi_j\}=2\delta_{ij}I,
\label{eq:SYKH}
\end{equation}
where $N$ is even, $2\leq q\leq N$ is even,
\begin{equation}
\mathcal I_N=\{(i_1,\ldots,i_q):1\le i_1<\cdots<i_q\le N\},
\qquad
g_I\overset{\mathrm{iid}}\sim\mathcal N(0,1),
\label{eq:variance}
\end{equation}
and
\begin{equation*}
\psi_I:=i^{q/2}\psi_{i_1}\cdots\psi_{i_q},
\qquad
\sigma_N^2=\frac{(q-1)!}{N^{q-1}}.
\end{equation*}
At finite $N$, this defines a random self-adjoint operator on a Hilbert space of dimension $2^{N/2}$.

The main question is how the thermal observables of this random matrix behave as $N$ tends to infinity. The matrix dimension grows exponentially, whereas the number of independent random coefficients grows polynomially for fixed $q$. The resulting correlations between matrix entries distinguish SYK from the usual ensembles of random matrix theory.

For $q=2$, a change of one-particle basis diagonalizes the Hamiltonian. For fixed $q\geq4$, the large-$N$ calculation instead leads to a nonlinear equation for a two-time correlation function. Its low-frequency approximation has a power-law solution. Fluctuations around that solution lead to the Schwarzian action, which also arises at the boundary of nearly-$\AdS_2$ gravity. Our purpose is to explain these reductions and the assumptions needed at each passage.

The organization is as follows. Section~\ref{sec:foundations} develops the finite-dimensional model. Section~\ref{sec:largeN} introduces the large-$N$ equations and their infrared approximation. Sections~\ref{sec:chaos} and~\ref{sec:gravity} discuss the physical interpretations. Section~\ref{sec:probability} returns to the probabilistic formulation and outstanding questions. The appendices supply algebraic calculations and a terminology guide. A first reading may stop after the Schwarzian action and resume with the probabilistic discussion.

Unless stated otherwise, $q$ is fixed and $N$ tends to infinity through even integers. Statements about the interacting infrared regime concern $q\geq4$, usually $q=4$. We use $\hbar=k_B=1$ and set the interaction scale to one. The symbols $\E$ and $\langle\cdot\rangle_{\beta,g}$ denote disorder expectation and thermal expectation at fixed disorder, respectively. We write $a_N\asymp b_N$ when their positive ratio stays between two constants independent of $N$.

\subsubsection{The mean-field analogy}

In the Curie-Weiss model, one starts from $N$ spins and obtains in the thermodynamic limit a variational problem for one scalar order parameter $m$. In the Sherrington-Kirkpatrick model, disorder introduces more elaborate order parameters, including overlap distributions. In SYK the large-$N$ order parameter is a two-time function $G(\tau_1,\tau_2)$. Thus the microscopic system is zero-dimensional in space, but the order parameter is infinite-dimensional because it depends on time.

The central reduction is
\begin{equation*}
\psi_1,\ldots,\psi_N
\quad\rightsquigarrow\quad
(G,\Sigma),
\end{equation*}
where $G$ is the Green kernel and $\Sigma$ is the self-energy. The collective action is proportional to $N$, so $1/N$ acts as a semiclassical fluctuation parameter. Strong coupling and analytic control coexist because the small parameter is not the interaction strength but the inverse number of fermion species.

\subsubsection{Status of the arguments}
Definitions, propositions, and their proofs concern ordinary finite-dimensional mathematics unless another setting is specified. A \emph{formal statement} records a conclusion of diagrammatic or saddle-point analysis without claiming a uniform error estimate. An \emph{effective description} concerns specified observables in a scaling regime; equality of effective actions alone does not identify microscopic Hilbert spaces. We use these distinctions throughout, especially when discussing the infrared limit and gravity.

\subsection{Fermions and Majorana operators}

We begin with exterior algebra, which gives a concrete representation of the fermionic operators.

\subsubsection{Antisymmetric Fock space}

Let the one-particle Hilbert space be $\mathfrak h=\C^M$. The fermionic Fock space is
\begin{equation}
\cF_-(\mathfrak h)
=\bigoplus_{r=0}^{M}\Lambda^r\mathfrak h.
\label{eq:fock}
\end{equation}
Its dimension is
\begin{equation*}
\dim\cF_-(\C^M)=\sum_{r=0}^{M}\binom Mr=2^M.
\end{equation*}
Choose an orthonormal basis $e_1,\ldots,e_M$. Exterior multiplication by $e_a$ defines the creation operator $c_a^\dagger\psi=e_a\wedge\psi$, and its adjoint $c_a$ is the annihilation operator.

\begin{proposition}[Canonical anticommutation relations]
The operators $c_a,c_a^\dagger$ satisfy
\begin{align*}
\{c_a,c_b^\dagger\}&=\delta_{ab}I,\\
\{c_a,c_b\}&=0,\\
\{c_a^\dagger,c_b^\dagger\}&=0.
\end{align*}
\end{proposition}

\begin{proof}
On a basis wedge, the adjoint of exterior multiplication is
\begin{equation*}
c_a(e_{j_1}\wedge\cdots\wedge e_{j_r})
=\sum_{\ell=1}^r(-1)^{\ell-1}\delta_{a j_\ell}
e_{j_1}\wedge\cdots\wedge\widehat e_{j_\ell}\wedge\cdots\wedge e_{j_r}.
\end{equation*}
Consequently $c_a(e_b\wedge u)=\delta_{ab}u-e_b\wedge c_a u$, which proves the first relation. Exterior multiplication anticommutes; taking adjoints gives the remaining relation.
\end{proof}

The nilpotence $(c_a^\dagger)^2=0$ is the algebraic form of Pauli exclusion. The number operator $n_a=c_a^\dagger c_a$ satisfies $n_a^2=n_a, \qquad \operatorname{spec}(n_a)=\{0,1\}$. A convenient orthonormal basis is $\ket{n_1,\ldots,n_M}, \qquad n_a\in\{0,1\}$. Thus $M$ complex fermionic modes have the same Hilbert-space dimension as $M$ qubits, but the natural generators satisfy anticommutation relations rather than commuting across tensor factors.

\subsubsection{Parity and the graded algebra}

The total number operator is
\begin{equation*}
\mathcal N=\sum_{a=1}^{M}n_a.
\end{equation*}
The fermion-parity operator is $(-1)^F=e^{i\pi\mathcal N}$. Operators which change particle number by an even amount commute with parity; odd operators anticommute with parity. Physical Hamiltonians are typically even. In SYK the quartic Hamiltonian is even, hence $[H,(-1)^F]=0$. This grading is important when discussing spectral degeneracies and random-matrix symmetry classes.

\subsubsection{Majorana fermions and Clifford algebras}

For each complex mode define two self-adjoint operators
\begin{align*}
\psi_{2a-1}&=c_a+c_a^\dagger,\\
\psi_{2a}&=i(c_a-c_a^\dagger).
\end{align*}
Then
\begin{equation}
\psi_i^*=\psi_i,
\qquad
\{\psi_i,\psi_j\}=2\delta_{ij}I,
\qquad
\psi_i^2=I.
\label{eq:majoranaCAR}
\end{equation}
Conversely,
\begin{equation*}
c_a=\frac{\psi_{2a-1}-i\psi_{2a}}2.
\end{equation*}
Hence $N=2M$ Majorana fermions act irreducibly on $\cH_N\cong\C^{2^{N/2}}$. Algebraically, they generate a complex Clifford algebra.

A product of $q$ distinct Majoranas has a simple adjoint rule. If $X=\psi_{i_1}\cdots\psi_{i_q}, \qquad i_1<\cdots<i_q$, then $X^*=(-1)^{q(q-1)/2}X$. For $q=4$, the sign is $+1$, so $\psi_i\psi_j\psi_k\psi_l$ is self-adjoint. For every even $q$, the phase in $\psi_I=i^{q/2}\psi_{i_1}\cdots\psi_{i_q}$ makes $\psi_I$ self-adjoint; moreover $\psi_I^2=I$ in the normalization used here.

\subsubsection{Jordan-Wigner representation}

For numerical intuition, one may represent Majoranas using Pauli matrices. On $(\C^2)^{\otimes M}$,
\begin{align*}
\psi_{2a-1}&=
\left(\prod_{b<a}\sigma_b^z\right)\sigma_a^x,\\
\psi_{2a}&=
\left(\prod_{b<a}\sigma_b^z\right)\sigma_a^y.
\end{align*}
The strings of $\sigma^z$ are precisely what generate anticommutation between different modes. This realizes the Clifford generators as concrete matrices on a tensor-product Hilbert space.

\subsubsection{Free fermions and quadratic Hamiltonians}

A general quadratic Majorana Hamiltonian can be written as
\begin{equation*}
H_2=\frac i4\sum_{i,j=1}^{N}A_{ij}\psi_i\psi_j,
\qquad
A^T=-A,
\qquad A\in M_N(\R).
\end{equation*}
This includes the $q=2$ SYK model: in that case the random couplings are simply the entries of the real skew-symmetric matrix $A$.  Its diagonalization follows from the canonical form of a real skew-symmetric matrix.  For $N=2M$, there are $O\in O(N)$ and numbers $\varepsilon_a\geq0$ such that
\begin{equation}
OAO^T=
\bigoplus_{a=1}^{M}
\begin{pmatrix}
0&-\varepsilon_a\\
\varepsilon_a&0
\end{pmatrix}.
\label{eq:skewCanonical}
\end{equation}
Thus the one-particle eigenvalues of $A$ are the pairs $\pm i\varepsilon_a$.  If
\begin{equation*}
\chi_i=\sum_jO_{ij}\psi_j,
\end{equation*}
then orthogonality gives $\{\chi_i,\chi_j\}=2\delta_{ij}I$, so this change of basis preserves the Majorana CAR.  Pairing the transformed Majoranas into complex fermions,
\begin{equation*}
d_a=\frac{\chi_{2a-1}-i\chi_{2a}}2,
\qquad n_a=d_a^\dagger d_a,
\end{equation*}
gives $i\chi_{2a-1}\chi_{2a}=1-2n_a$ and hence
\begin{equation}
H_2=\sum_{a=1}^{M}\varepsilon_a
\left(n_a-\frac12\right).
\label{eq:freeFermionDiagonal}
\end{equation}
In particular, $[n_a,n_b]=0$.  Every many-body eigenvector is an occupation-number state $\ket{n_1,\ldots,n_M}, \qquad n_a\in\{0,1\}$, with the additive energy
\begin{equation*}
E(n_1,\ldots,n_M)
=\sum_{a=1}^{M}\varepsilon_a
\left(n_a-\frac12\right).
\end{equation*}
Therefore diagonalizing the exponentially large $2^M\times2^M$ many-body Hamiltonian reduces to diagonalizing the $N\times N$ matrix $A$.  For example, the partition function factorizes as
\begin{equation*}
Z_2(\beta)=\prod_{a=1}^{M}
2\cosh\left(\frac{\beta\varepsilon_a}{2}\right),
\end{equation*}
and all thermal correlation functions follow from the two-point function by Wick's theorem.  This is the precise sense in which $q=2$ is a free-fermion, and hence explicitly solvable, model.

The same simplicity is visible dynamically.  The commutator of a quadratic Hamiltonian with a Majorana operator is again linear:
\begin{equation*}
[H_2,\psi_k]=-i\sum_jA_{kj}\psi_j.
\end{equation*}
Consequently the Heisenberg equations close on the $N$-dimensional one-particle space,
\begin{equation*}
\frac{d}{dt}\psi(t)=A\psi(t),
\qquad
\psi(t)=e^{tA}\psi(0).
\end{equation*}
No new many-body operators are generated by time evolution.

For interacting SYK, the same argument no longer applies.  If $q\geq4$, then $[H_q,\psi_m]$ is generically a linear combination of products of $q-1$ Majoranas, rather than a linear combination of the $\psi_j$.  An orthogonal change of one-particle basis preserves polynomial degree: it can put a two-form $A\in\Lambda^2\R^N$ into the block form \eqref{eq:skewCanonical}, but it cannot turn a generic tensor in $\Lambda^q\R^N$ into a quadratic Hamiltonian.  The Heisenberg evolution therefore mixes increasingly complicated operator sectors, and the reduction to independent occupation numbers fails.  This is why the strongly interacting low-energy state need not admit a quasiparticle interpretation.

\subsection[Grassmann algebra and correlators]{Grassmann variables and thermal correlators}

The operator formulation is fundamental. Grassmann variables are a powerful algebraic device for reorganizing traces and correlation functions.

\subsubsection{Grassmann algebra and Berezin integration}

Let $\theta_1,\ldots,\theta_n$ generate a Grassmann algebra: $\theta_i\theta_j=-\theta_j\theta_i$. In particular, $\theta_i^2=0$. Berezin integration is defined by
\begin{equation*}
\int d\theta\,1=0,
\qquad
\int d\theta\,\theta=1.
\end{equation*}
It is differentiation and integration at the same time because every polynomial is affine in each Grassmann generator.

For an antisymmetric $2m\times2m$ matrix $A$,
\begin{equation}
\int d\theta_{2m}\cdots d\theta_1
\exp\left(\frac12\theta^T A\theta\right)
=\Pf(A),
\label{eq:pfaff}
\end{equation}
with $\Pf(A)^2=\det A$. For comparison, if $A$ is a real symmetric positive-definite $n\times n$ matrix, the ordinary Gaussian integral is
\begin{equation*}
\int_{\R^n}e^{-x^TAx/2}\,dx
=(2\pi)^{n/2}(\det A)^{-1/2}.
\end{equation*}
The determinant occurs with the opposite power for fermions. This algebraic fact is the source of the logarithmic determinant in the large-$N$ collective action.

\subsubsection{Finite-temperature quantum mechanics}

For a fixed disorder realization $g=(g_I)_{I\in\mathcal I_N}$, define
\begin{equation*}
Z_g(\beta)=\Tr e^{-\beta H_q},
\qquad
\rho_{\beta,g}=\frac{e^{-\beta H_q}}{Z_g(\beta)}.
\end{equation*}
For an observable $A$, $\langle A\rangle_{\beta,g}=\Tr(\rho_{\beta,g}A)$. Imaginary-time evolution is $A(\tau)=e^{\tau H}Ae^{-\tau H}, \qquad 0\le \tau<\beta$. The Euclidean two-point function is
\begin{equation*}
G_{ij,g}(\tau)=
\langle T_\tau \psi_i(\tau)\psi_j(0)\rangle_{\beta,g}.
\end{equation*}
After disorder averaging and flavor averaging,
\begin{equation}
G_N(\tau)=\frac1N\sum_{i=1}^{N}\E_g G_{ii,g}(\tau).
\label{eq:GN}
\end{equation}
Fermionic thermal correlators are antiperiodic, $G_N(\tau+\beta)=-G_N(\tau)$, which implies the Matsubara frequencies
\begin{equation*}
\omega_n=\frac{(2n+1)\pi}{\beta}.
\end{equation*}
Thus
\begin{equation*}
G(\tau)=\frac1\beta\sum_{n\in\mathbb Z}e^{-i\omega_n\tau}G(i\omega_n).
\end{equation*}

\subsubsection{Spectral representation}

Fix the disorder and a self-adjoint Majorana operator $\psi$. Let $H\ket n=E_n\ket n$ and $Z=\sum_n e^{-\beta E_n}$. Its positive spectral measure is
\begin{equation}
d\nu_\psi(\omega)=\frac1Z\sum_{m,n}
(e^{-\beta E_n}+e^{-\beta E_m})
|\langle n|\psi|m\rangle|^2\,
\delta_{E_m-E_n}(d\omega).
\label{eq:spectralmeasure}
\end{equation}
In particular, $\nu_\psi(\R)=\langle\{\psi,\psi\}\rangle_\beta=2$.
With our convention $G(\tau)=\langle T_\tau\psi(\tau)\psi(0)\rangle_\beta$ and Fourier transform $G(i\omega_n)=\int_0^\beta e^{i\omega_n\tau}G(\tau)\,d\tau$, the spectral formulas are
\begin{equation}
G(i\omega_n)=-\int_\R\frac{d\nu_\psi(\omega)}{i\omega_n-\omega},
\qquad
G^R(z)=\int_\R\frac{d\nu_\psi(\omega)}{z-\omega},\quad \operatorname{Im}z>0.
\label{eq:spectral}
\end{equation}
Here $G^R(t)=-i\theta(t)\langle\{\psi(t),\psi(0)\}\rangle_\beta$. Thus the retarded function is the continuation of $-G$, rather than $G$, in this Euclidean sign convention. If the spectral measure has a density $A$, then $A(\omega)=-\pi^{-1}\operatorname{Im}G^R(\omega+i0^+)$. At finite $N$, this identity is understood distributionally because the measure is atomic. Appendix~\ref{app:identities} gives the corresponding imaginary-time expansion.

\subsection{The SYK Hamiltonian}

\subsubsection{Definition and normalization}

Fix even integers $N$ and $q$ with $2\leq q\leq N$, and let $\mathcal I_N=\{(i_1,\ldots,i_q):1\le i_1<\cdots<i_q\le N\}$. For $I=(i_1,\ldots,i_q)\in\mathcal I_N$, set $\psi_I:=i^{q/2}\psi_{i_1}\cdots\psi_{i_q}$. Let $(g_I)_{I\in\mathcal I_N}$ be iid standard normal random variables and define
\begin{equation}
H_q=\sigma_N\sum_{I\in\mathcal I_N}g_I\psi_I,
\qquad
\sigma_N^2=\frac{(q-1)!}{N^{q-1}}.
\label{eq:Hqnormalized}
\end{equation}
Equivalently, the random coefficients $J_I:=\sigma_Ng_I$ are independent centered Gaussians with
\begin{equation}
\E[J_IJ_{I'}]
=\frac{(q-1)!}{N^{q-1}}\delta_{I,I'}.
\label{eq:Jcovariance}
\end{equation}
Because $q$ is even and $\{\psi_i,\psi_j\}=2\delta_{ij}I$, each $\psi_I$ is a self-adjoint unitary. For $q=4$, formula \eqref{eq:Hqnormalized} differs by an irrelevant overall sign from writing the interaction as a sum of $\psi_i\psi_j\psi_k\psi_l$, since the centered Gaussian law is invariant under $g_I\mapsto-g_I$.

\begin{example}[Two Majorana operators]
For $N=q=2$, take $\psi_1=\sigma^x$ and $\psi_2=\sigma^y$. Since $i\psi_1\psi_2=-\sigma^z$, the Hamiltonian and partition function are
\begin{equation*}
H_2=-\frac{g}{\sqrt2}\sigma^z,
\qquad
Z_g(\beta)=2\cosh\left(\frac{\beta g}{\sqrt2}\right),
\qquad g\sim\mathcal N(0,1).
\end{equation*}
Thus thermal expectation is a trace over two states, while disorder expectation is an integral against one Gaussian variable. These are the two averages that remain distinct in the general model.
\end{example}

\subsubsection{Why the variance scales as $N^{-(q-1)}$}

There are
\begin{equation*}
\binom Nq\sim \frac{N^q}{q!}
\end{equation*}
interaction monomials. The variance scaling is chosen so that the large-$N$ self-energy remains of order one and the total free energy is of order $N$. In diagrammatic language, each disorder contraction contributes $N^{-(q-1)}$, while free flavor-index sums contribute powers of $N$. The competition produces a nontrivial $N\to\infty$ limit rather than either a free theory or a divergent one.

A second-moment computation makes the scale precise. Since $\psi_I^2=I$ and distinct Clifford monomials are orthogonal under the normalized trace,
\begin{equation*}
2^{-N/2}\E\Tr H_q^2
=\sigma_N^2\binom Nq
=\frac Nq\prod_{r=1}^{q-1}\left(1-\frac rN\right)
\sim\frac Nq.
\end{equation*}
The root-mean-square eigenvalue, averaged over disorder and normalized trace, is therefore of order $\sqrt N$. This calculation does not determine the spectral edge. The order-$N$ ground-state scale requires a separate argument; see Section~\ref{sec:probability}.

\subsubsection{Quenched and annealed thermodynamics}

The quenched free-energy density is
\begin{equation*}
f_N^{\mathrm q}(\beta)
=-\frac1{\beta N}\E\log Z_g(\beta).
\end{equation*}
The annealed free-energy density is
\begin{equation*}
f_N^{\mathrm a}(\beta)
=-\frac1{\beta N}\log\E Z_g(\beta).
\end{equation*}
By Jensen, $\E\log Z\le \log\E Z$, so $f_N^{\mathrm a}(\beta)\le f_N^{\mathrm q}(\beta)$. Concentration of the sample pressure is elementary in the Gaussian model. It does not, by itself, imply equality of the quenched and annealed limits.

\begin{proposition}[Concentration of the pressure]\label{prop:concentration}
For $\beta>0$, put $P_N(g)=N^{-1}\log Z_g(\beta)$. Then
\begin{equation*}
\operatorname{Var}(P_N)\leq\frac{\beta^2}{qN},
\qquad
\Pp\bigl(|P_N-\E P_N|\geq t\bigr)
\leq 2\exp\left(-\frac{qNt^2}{2\beta^2}\right),\quad t>0.
\end{equation*}
In particular, $P_N-\E P_N\to0$ in probability at every fixed $\beta$.
\end{proposition}

\begin{proof}
Differentiation of the matrix exponential under the trace, followed by cyclicity, gives
\begin{equation*}
\partial_{g_I}P_N=-\frac{\beta\sigma_N}{N}\langle\psi_I\rangle_{\beta,g}.
\end{equation*}
Since $\|\psi_I\|=1$,
\begin{equation*}
\|\nabla P_N\|_2^2
\leq\frac{\beta^2\sigma_N^2}{N^2}\binom Nq
=\frac{\beta^2}{qN}\prod_{r=1}^{q-1}\left(1-\frac rN\right)
\leq\frac{\beta^2}{qN}.
\end{equation*}
The Gaussian Poincar\'e inequality proves the variance bound, and Gaussian concentration for Lipschitz functions proves the tail bound.
\end{proof}

\begin{remark}[Dependence on temperature]
The pressure bound in Proposition~\ref{prop:concentration} gives concentration along every sequence $\beta_N=o(\sqrt N)$. This is a sufficient range for that bound, not a physical crossover. Indeed, the sample free-energy density $f_{N,g}(\beta)=-P_N(g)/\beta$ satisfies
\begin{equation*}
\operatorname{Var}(f_{N,g}(\beta))\leq\frac1{qN},
\qquad
\Pp\bigl(|f_{N,g}(\beta)-\E f_{N,g}(\beta)|\geq t\bigr)
\leq2e^{-qNt^2/2}
\end{equation*}
for every $\beta>0$. Thus its concentration persists for arbitrary deterministic temperature sequences. The loss of a useful estimate for one observable does not imply a loss of self-averaging for all observables, nor does it locate the Schwarzian crossover.
\end{remark}

Macroscopic self-averaging is compatible with microscopic random-matrix fluctuations. The free energy can concentrate while individual level spacings remain random.

\section[Large-N and infrared theory]{The large-$N$ equations and the infrared limit}\label{sec:largeN}

\subsection[The melonic large-N limit]{The melonic large-$N$ limit}

The key solvability mechanism is a combinatorial selection rule on Feynman graphs.

\subsubsection{Disorder averaging}

In a perturbative expansion, each interaction vertex carries a random coefficient $\sigma_Ng_I$. Since the disorder is Gaussian, averaging pairs the couplings: $\E[(\sigma_Ng_I)(\sigma_Ng_{I'})] =\sigma_N^2\delta_{I,I'}$. After averaging, every pair of interaction vertices comes with identified flavor indices. The resulting graph carries a power of $N$ determined by how many independent flavor-index loops remain.

The leading graphs are the melonic graphs. They are recursively generated by replacing a propagator line with a two-vertex self-energy insertion and iterating.

\begin{heuristic}[Melonic dominance]
At fixed $q$ and fixed perturbative order, the leading power of $N$ is attained by melonic graphs. Nonmelonic graphs are suppressed by powers of $1/N$.
\end{heuristic}

For SYK, melonic dominance has a rigorous combinatorial diagrammatic proof~\cite{Bonzom}. What remains formal is the passage from order-by-order diagrammatics to the full nonperturbative infrared theory in the strongest sense used in physics.

\subsubsection{Dyson equation}

Under $\{\psi_i,\psi_j\}=2\delta_{ij}I$, one has $G_0(0^+)=1$ and the free Majorana propagator is $G_0(i\omega_n)=2(-i\omega_n)^{-1}$. The exact Dyson relation is $G(i\omega_n)^{-1}=G_0(i\omega_n)^{-1}-\Sigma(i\omega_n)$. With the convention used here,
\begin{equation}
G(i\omega_n)^{-1}=-\frac{i\omega_n}{2}-\Sigma(i\omega_n).
\label{eq:dyson}
\end{equation}
Melonic closure gives for quartic SYK
\begin{equation}
\Sigma(\tau)=G(\tau)^3.
\label{eq:selfenergy}
\end{equation}
For general even $q$, $\Sigma(\tau)=G(\tau)^{q-1}$. The unit coefficient follows from using both $\sigma_N^2=(q-1)!/N^{q-1}$ and Clifford generators squaring to $I$.

Equations \eqref{eq:dyson} and \eqref{eq:selfenergy} are the central large-$N$ Schwinger-Dyson equations.

\subsubsection{Kernel form}

Regard $G$ and $\Sigma$ as integral kernels. Then $(\tfrac12\partial_\tau-\Sigma)\circ G=I$, meaning
\begin{equation*}
\int_0^\beta d\tau'
\left[
\frac12\partial_{\tau_1}\delta(\tau_1-\tau')-\Sigma(\tau_1,\tau')
\right]G(\tau',\tau_2)
=\delta(\tau_1-\tau_2).
\end{equation*}
At a translation-invariant thermal saddle, $G(\tau_1,\tau_2)=G(\tau_1-\tau_2)$. Thus a many-body problem on a Hilbert space of dimension $2^{N/2}$ becomes a nonlinear integral equation for one scalar function of time.

\subsection{The bilocal collective action}

There is a second derivation of the same equations which is closer to mean-field variational methods.

\subsubsection{Introducing $G$ and $\Sigma$}

Formally insert the constraint
\begin{equation*}
G(\tau_1,\tau_2)
=\frac1N\sum_{i=1}^{N}\psi_i(\tau_1)\psi_i(\tau_2)
\end{equation*}
using a bilocal Lagrange multiplier $\Sigma(\tau_1,\tau_2)$. After Gaussian disorder averaging and integrating out the fermions, one obtains for general even $q$
\begin{equation}
\frac{I_{\mathrm{eff}}[G,\Sigma]}{N}
=-\frac12\log\det(\tfrac12\partial_\tau-\Sigma)
+\frac12\int d\tau_1d\tau_2
\left[
\Sigma G-\frac1qG^q
\right].
\label{eq:collective}
\end{equation}
The determinant term is the logarithm of a Pfaffian in Majorana language. Equation \eqref{eq:pfaff} is the finite-dimensional algebraic prototype.

For a single copy, the disorder average gives the annealed partition function in the formal form
\begin{equation*}
\E Z\simeq\int \mathcal DG\,\mathcal D\Sigma\,
\exp\left[-N\mathcal I(G,\Sigma)\right].
\end{equation*}
Here $\mathcal I=I_{\mathrm{eff}}/N$, and a regularization and integration contour are implicit. Quenched observables require replicas or a separate justification that the relevant saddle gives their limit. See~\cite{MaldacenaStanford,KitaevSuh} for the collective-field calculation.

\subsubsection{Stationary equations}

Varying \eqref{eq:collective} with respect to $G$ gives $\Sigma(\tau_1,\tau_2)=G(\tau_1,\tau_2)^{q-1}$. Varying with respect to $\Sigma$ gives $G=(\tfrac12\partial_\tau-\Sigma)^{-1}$. Hence the saddle equations are precisely the Schwinger-Dyson equations.

\begin{remark}[Analogy with a Laplace principle]
The schematic similarity
\begin{equation*}
\int e^{-NI(x)}dx\sim e^{-N\inf I}
\end{equation*}
is useful, but the fermionic functional integral is not automatically a rigorous probability measure. A mathematical derivation requires control of disorder averaging, the collective-variable transformation, saddle selection, and low-temperature limits.
\end{remark}

\subsubsection{Large-$N$ factorization}

One consequence of the prefactor $N$ is that fluctuations of suitably normalized singlet observables are suppressed. Schematically,
\begin{equation*}
\langle \cO_1\cO_2\rangle
-\langle\cO_1\rangle\langle\cO_2\rangle
=O(N^{-1}).
\end{equation*}
This is the analogue of concentration around a mean-field saddle. The connected four-point function is therefore naturally an $O(1/N)$ object.

\subsection{The infrared conformal limit}

We now take $q\geq4$ and examine the formal low-frequency approximation to the Schwinger--Dyson equations. The approximation is made after the large-$N$ reduction~\cite{MaldacenaStanford}.

\subsubsection{Dropping the kinetic term}

The full equation is
\begin{equation*}
G(i\omega)^{-1}=-\frac{i\omega}{2}-\Sigma(i\omega).
\end{equation*}
If $|\omega|\ll1$, and if the self-energy dominates $i\omega$, one may first solve the scale-invariant equation obtained by dropping the kinetic term. In time-domain kernel notation,
\begin{equation}
\Sigma\circ G=-I.
\label{eq:IRconv}
\end{equation}
Together with $\Sigma(\tau)=G(\tau)^{q-1}$, this admits a power-law solution.

Assume
\begin{equation*}
G_c(\tau)=b\,\frac{\sgn\tau}{|\tau|^{2\Delta}}.
\end{equation*}
Then
\begin{equation*}
\Sigma_c(\tau)\propto \frac{\sgn\tau}{|\tau|^{2\Delta(q-1)}}.
\end{equation*}
The convolution in \eqref{eq:IRconv} must scale as a delta distribution. Since $\delta(\lambda\tau)=\lambda^{-1}\delta(\tau)$, dimensional consistency gives $2\Delta(q-1)+2\Delta=2$, so
\begin{equation*}
\Delta=\frac1q.
\end{equation*}
For $q=4$,
\begin{equation}
\Delta=\frac14,
\qquad
G_c(\tau)=b\frac{\sgn\tau}{|\tau|^{1/2}}.
\label{eq:IRG}
\end{equation}
With the normalization used here,
\begin{equation*}
b^4=\frac1{4\pi}.
\end{equation*}

\subsubsection{Finite temperature}

The conformal map from the line to the thermal circle suggests
\begin{equation}
G_{\beta,c}(\tau)
=b\,\sgn(\tau)
\left[
\frac{\pi}{\beta\sin(\pi|\tau|/\beta)}
\right]^{2\Delta},
\qquad 0<|\tau|<\beta.
\label{eq:finiteTG}
\end{equation}
For $q=4$, the exponent is $1/2$.

The physical interpretation is that strong interactions generate an anomalous scaling dimension. The fermion field does not behave like a weakly dressed free particle.

\subsubsection{Emergent reparametrization covariance}

The infrared equations are covariant under orientation-preserving time reparametrizations $\tau\mapsto f(\tau)$ if
\begin{equation}
G(\tau_1,\tau_2)
\mapsto
[f'(\tau_1)f'(\tau_2)]^\Delta
G(f(\tau_1),f(\tau_2)),
\label{eq:reparamG}
\end{equation}
while
\begin{equation*}
\Sigma(\tau_1,\tau_2)
\mapsto
[f'(\tau_1)f'(\tau_2)]^{1-\Delta}
\Sigma(f(\tau_1),f(\tau_2)).
\end{equation*}
A particular conformal saddle is invariant only under the M\"obius subgroup
\begin{equation*}
f(\tau)=\frac{a\tau+b}{c\tau+d},
\qquad ad-bc=1.
\end{equation*}
Thus the symmetry pattern is often written
\begin{equation*}
\Diff(S^1)\longrightarrow PSL(2,\R).
\end{equation*}
The full microscopic theory is not exactly reparametrization invariant. The kinetic term explicitly breaks the symmetry. Therefore the reparametrization modes are soft rather than exact zero modes.

\subsection{The Schwarzian soft mode}\label{sec:soft}

\subsubsection{The Schwarzian derivative}

For a smooth monotone function $f$ define
\begin{equation}
\{f,\tau\}
=\frac{f'''(\tau)}{f'(\tau)}
-\frac32\left(\frac{f''(\tau)}{f'(\tau)}\right)^2.
\label{eq:schwarzian}
\end{equation}
It satisfies the projective invariance
\begin{equation*}
\left\{\frac{af+b}{cf+d},\tau\right\}=\{f,\tau\}.
\end{equation*}
Hence $PSL(2,\R)$ is naturally quotiented out.

The leading low-energy action for reparametrization fluctuations has the form
\begin{equation}
I_{\Sch}[f]
=-N\alpha_S
\int_0^\beta d\tau
\left\{
\tan\left(\frac{\pi f(\tau)}{\beta}\right),\tau
\right\},
\label{eq:SchAction}
\end{equation}
Here $\alpha_S>0$ is a model-dependent constant, and we write $C=N\alpha_S$ in the units of these notes. For orientation-preserving maps with $f(\tau+\beta)=f(\tau)+\beta$, the formal configuration space is $\Diff(S^1)/PSL(2,\R)$.

\subsubsection{Why the Schwarzian matters}

The strict conformal approximation produces a nearly flat manifold of reparametrized saddles. The small UV kinetic term lifts this degeneracy. The most infrared-singular fluctuations therefore come from motion along this almost-flat manifold, not from generic massive fluctuations. Integrating out the massive directions produces the Schwarzian action.

\subsubsection{The dimensionless fluctuation parameter}
Put $u=\tau/\beta$ and $f(\tau)=\beta g(u)$, where $g(u+1)=g(u)+1$. The chain rule for the Schwarzian gives
\begin{equation}
I_{\Sch}[f]=\frac C\beta\,\mathcal A[g],
\qquad
\mathcal A[g]:=-\int_0^1\{\tan(\pi g(u)),u\}\,du.
\label{eq:SchRescaled}
\end{equation}
This identifies $\beta/C$ as the fluctuation parameter. For a fixed smooth periodic function $h$, set $g(u)=u+\varepsilon h(u)$. Integration by parts gives
\begin{align*}
\mathcal A[g]
&=\frac12\int_0^1\left(\frac{g''}{g'}\right)^2du
-2\pi^2\int_0^1(g')^2du,\\
\mathcal A[u+\varepsilon h]
&=-2\pi^2+\frac{\varepsilon^2}{2}
\int_0^1\bigl((h'')^2-4\pi^2(h')^2\bigr)\,du
+O_h(\varepsilon^3).
\end{align*}
The quadratic form vanishes on the Fourier modes with indices $0,1,-1$, which generate the projective redundancy. On the remaining modes it is positive. The Gaussian approximation therefore assigns dimensionless fluctuations of size $(\beta/C)^{1/2}$ to a fixed nonzero mode after this redundancy is removed.

For $\beta/C\to0$, expansion around one thermal saddle is appropriate. When $\beta/C$ is of order one, the soft coordinate must be integrated over. This is a singular saddle-point problem: the leading collective action has a family of stationary points, and its small explicit symmetry breaking determines an action along that family. The coordinate $f$ parametrizes collective configurations; it is not an additional microscopic fermion.

\subsubsection{Low-temperature thermodynamics}

At leading order in $N$, in the semiclassical range $1\ll\beta\ll N$, the Schwarzian theory gives $F(T)=E_0-S_0T-cT^2+\cdots$, so
\begin{equation*}
S(T)=-\frac{\partial F}{\partial T}
=S_0+2cT+\cdots.
\end{equation*}
For SYK at large $N$,
\begin{equation}
S(T)=Ns_0+N\gamma T+\cdots.
\label{eq:SYKentropy}
\end{equation}
The residual entropy density $s_0$ is nonzero in the order of limits
\begin{equation*}
\lim_{T\downarrow0}\lim_{N\to\infty}\frac{S_N(T)}N=s_0>0.
\end{equation*}
For generic nonsupersymmetric finite-$N$ SYK, exact ground-state degeneracy is not exponentially large. Therefore
\begin{equation*}
\lim_{N\to\infty}\lim_{T\downarrow0}\frac{S_N(T)}N=0
\end{equation*}
in the generic situation, up to symmetry degeneracies. The two limits do not commute.

\subsubsection{The quantum Schwarzian partition function}
The Schwarzian functional integral can be evaluated beyond the small-fluctuation regime~\cite{StanfordWitten}. Choose its overall normalization so that
\begin{equation}
Z_{\Sch}(\beta;C)
=\frac1{\sqrt{2\pi}}\left(\frac C\beta\right)^{3/2}
e^{2\pi^2C/\beta}.
\label{eq:SchPartition}
\end{equation}
Its Laplace density is then
\begin{equation}
\rho_{\Sch}(\varepsilon;C)
=\frac{C}{2\pi^2}\sinh\bigl(2\pi\sqrt{2C\varepsilon}\bigr),
\qquad \varepsilon\geq0.
\label{eq:SchDensity}
\end{equation}
Indeed, for $a,\beta>0$,
\begin{equation*}
\int_0^\infty e^{-\beta\varepsilon}\sinh(a\sqrt\varepsilon)\,d\varepsilon
=\frac{a\sqrt\pi}{2\beta^{3/2}}e^{a^2/(4\beta)},
\end{equation*}
as follows by differentiating the Gaussian integral after setting $\varepsilon=x^2$. Taking $a=2\pi\sqrt{2C}$ verifies the normalization.

Let $E_0$ and $S_0=Ns_0(q)+o(N)$ be the smooth edge and entropy parameters. Subextensive terms in $S_0$ depend on normalization. The effective partition function is
\begin{equation}
Z_{\mathrm{smooth}}(\beta)
=e^{S_0-\beta E_0}Z_{\Sch}(\beta;C),
\qquad
\rho_{\mathrm{smooth}}(E_0+\varepsilon)
=e^{S_0}\rho_{\Sch}(\varepsilon;C).
\label{eq:SchMatching}
\end{equation}
Here $S_0$ is already extensive. A different overall normalization of $Z_{\Sch}$ must be accompanied by the corresponding change in $S_0$. In particular,
\begin{align*}
\log Z_{\mathrm{smooth}}
&=S_0-\beta E_0+\frac{2\pi^2C}{\beta}
+\frac32\log\frac C\beta-\frac12\log(2\pi),\\
-\partial_\beta\log Z_{\mathrm{smooth}}
&=E_0+\frac{2\pi^2C}{\beta^2}+\frac{3}{2\beta}.
\end{align*}
The term $\tfrac32\log C$ cannot be absorbed into a constant independent of $N$. At leading semiclassical order these formulas give $c=2\pi^2C$ and $\gamma=4\pi^2\alpha_S$ in \eqref{eq:SYKentropy}.

Formula~\eqref{eq:SchPartition} is exact within the Schwarzian theory, with its specified measure. It is an approximation to microscopic SYK only in a regime where matching to that theory is justified. Its smooth Laplace density also differs from the atomic spectral measure of a finite matrix. For $\beta\gg C$, the Schwarzian coordinate fluctuates strongly, while a smooth spectral description can still be useful if the relevant energy window contains many levels.

\subsubsection{A hierarchy of large-$N$ and low-temperature regimes}

The phrases ``large $N$'' and ``low temperature'' do not specify a unique asymptotic regime.  Their meaning depends on how the inverse temperature $\beta$ grows with $N$, and the corresponding changes should not all be called phase transitions.  To restore an interaction scale, replace $H_q$ by $JH_q$ with $J>0$. Then $C=N\alpha_S/J$, and the dimensionless inverse temperature is $\beta J$. We otherwise retain $J=1$ and the convention $\psi_i^2=I$.  The principal ranges, with the last crossover estimated from the smooth density, are
\begin{center}
\begin{tabular}{@{}ll@{}}
\toprule
inverse-temperature scale & effective description \\
\midrule
$\beta J=O(1)$ & ordinary large-$N$ Schwinger--Dyson theory \\
$1\ll\beta J\ll N$ & conformal SYK and a semiclassical Schwarzian mode \\
$\beta J\asymp N$ & Schwarzian fluctuations become order one \\
$N\ll\beta J\ll N e^{2S_0/3}$
    & quantum Schwarzian with an effectively continuous spectrum \\
$\beta J\asymp N e^{2S_0/3}$
    & the microscopic spectral edge becomes visible \\
$\beta J\gg N e^{2S_0/3}$
    & the lowest individual eigenvalues dominate \\
\bottomrule
\end{tabular}
\end{center}
Here $S_0$ is as in \eqref{eq:SchMatching}.  For the standard quartic Majorana model, $s_0(4)\simeq0.2324$.  The scales $\beta J\asymp1$ and $\beta J\asymp N$ describe crossovers between asymptotic approximations, not thermodynamic phase transitions.  The first is the UV-to-IR crossover.  At the second, the Schwarzian coupling $C\asymp N/J$ satisfies $\beta/C=O(1)$, so the reparametrization mode must be treated quantum mechanically rather than semiclassically.

The exponentially large scale in the table is a prediction obtained by extrapolating a smooth spectral density. It is not a proved local-spacing law for finite-$N$ SYK. The parameter $E_0$ in \eqref{eq:SchMatching} is not a sample ground-state eigenvalue.  The Schwarzian density near the lower spectral edge $E_0$ has the square-root form
\begin{equation}
\rho(E_0+\varepsilon)
\asymp e^{S_0}C^{3/2}\varepsilon^{1/2},
\qquad \varepsilon\downarrow0,
\label{eq:introSchEdge}
\end{equation}
up to constants independent of $\varepsilon$.  Consequently, the smooth density predicts
\begin{equation*}
\int_0^\varepsilon\rho(E_0+u)\,du
\asymp e^{S_0}C^{3/2}\varepsilon^{3/2}.
\end{equation*}
The continuum approximation reaches the first-level scale when this quantity is of order one.  Thus
\begin{equation}
\varepsilon_*
\asymp C^{-1}e^{-2S_0/3},
\qquad
\beta_{\mathrm{edge}}
\asymp C e^{2S_0/3},
\label{eq:introEdgeScale}
\end{equation}
At the level of this effective-theory estimate,
\begin{equation}
\lim_{N\to\infty}\frac1N
\log(\beta_{\mathrm{edge}}J)=\frac23s_0(q).
\label{eq:introEdgeRate}
\end{equation}
For $q=4$ the exponential rate is approximately $0.1549$.  Polynomial prefactors depend on conventions and on the observable used to locate the crossover; the exponent $2s_0/3$ follows from the assumed square-root edge.  It combines the exponential number $e^{S_0}$ of low-energy states with the square-root spectral edge.  More generally, an edge law $\rho(E_0+\varepsilon)\asymp e^{S_0}\varepsilon^\alpha$ would give the exponential rate $s_0/(\alpha+1)$.

The order-of-limits distinction in \eqref{eq:SYKentropy} is therefore a statement about spectral resolution. At fixed $\beta$, the large-$N$ limit cannot distinguish individual low-lying eigenvalues. Within the smooth approximation, the number of levels in the quantum-Schwarzian thermal window is of order
\begin{equation*}
\int_0^{1/\beta}\rho_{\mathrm{smooth}}(E_0+\varepsilon)\,d\varepsilon
\asymp e^{S_0}\left(\frac C\beta\right)^{3/2},
\qquad C\ll\beta.
\end{equation*}
It remains large when $C\ll\beta\ll\beta_{\mathrm{edge}}$. Thus order-one fluctuations of the soft mode occur at a polynomial scale, whereas individual edge levels are predicted to become visible at an exponential scale.

An assumption such as $\liminf N^{-1}\log(\beta_NJ)>0$ alone does not locate the discrete-spectrum regime. Its exponential rate must be compared with $2s_0(q)/3$, and its prefactor matters at equality. For a fixed sample, write the distinct energies as $\mathcal E_0<\mathcal E_1<\cdots$, with multiplicities $d_0,d_1,\ldots$. The exact identity
\begin{equation*}
Z_g(\beta)=d_0e^{-\beta\mathcal E_0}
\left(1+\sum_{a\geq1}\frac{d_a}{d_0}
e^{-\beta(\mathcal E_a-\mathcal E_0)}\right)
\end{equation*}
shows that ground-state dominance requires the remaining sum to be small. In a joint limit, even this condition needs control of both gaps and multiplicities. The smooth density predicts a candidate scale; it does not prove that condition for typical disorder. Annealed quantities such as $\log\E Z$ require additional control of rare samples~\cite{Cotler,SSS}.

\subsection{The four-point ladder kernel}

The two-point function determines the mean field. The four-point function probes interactions between collective excitations and is the gateway to the operator spectrum and chaos.

\subsubsection{Connected four-point function}

For distinct flavor indices in an averaged singlet channel, a normalized connected four-point function has size $\cF=O(N^{-1})$. Its leading diagrams form a ladder series. If $K$ is the integral operator adding one rung, then
\begin{equation}
\cF=\cF_0+K\cF_0+K^2\cF_0+\cdots
=(I-K)^{-1}\cF_0,
\label{eq:laddergeom}
\end{equation}
whenever the series is interpreted in a suitable spectral sense.

For quartic Majorana SYK, the conformal kernel can be written schematically as
\begin{equation*}
K(\tau_1,\tau_2;\tau_3,\tau_4)
=-3
G(\tau_{13})G(\tau_{24})G(\tau_{34})^2,
\end{equation*}
with convention-dependent signs.

\subsubsection{Conformal diagonalization}

In the conformal limit, $K$ commutes with the $SL(2,\R)$ action. Hence it can be diagonalized using one-dimensional conformal eigenfunctions. For the $q=4$ model, a commonly used eigenvalue function is
\begin{equation}
k(h)
=-\frac32
\frac{\tan\left(\frac\pi2(h-\frac12)\right)}{h-\frac12}.
\label{eq:kh}
\end{equation}
The dimensions $h$ of fermion bilinear operators satisfy
\begin{equation}
k(h)=1.
\label{eq:kh1}
\end{equation}
The special solution $h=2$ corresponds to the reparametrization soft mode. The strict conformal kernel has eigenvalue one on this sector. The kinetic correction lifts the degeneracy. Inverting $I-K$ enhances the corresponding response: for thermal-time separations, normalized soft contributions can scale as $\beta/C\asymp\beta J/N$, rather than uniformly as $1/N$. This is the same fluctuation parameter as in \eqref{eq:SchRescaled}.

\subsubsection{Interpretation as an operator-spectrum problem}

In a conformal theory one expects an operator product expansion of the schematic form
\begin{equation*}
\psi_i(\tau_1)\psi_i(\tau_2)
\sim
\sum_n C_n |\tau_{12}|^{h_n-2\Delta}\cO_{h_n}(\tau_2).
\end{equation*}
The transcendental equation \eqref{eq:kh1} determines the tower $h_n$. Thus a nontrivial interacting spectrum is reduced to the spectral analysis of an explicit integral operator.

This reduces the formal calculation of bilinear dimensions to an integral-operator eigenvalue problem~\cite{PolchinskiRosenhaus}. Its domain and the treatment of the $h=2$ sector must be specified before interpreting the resolvent analytically.

\section[Chaos and quasiparticles]{Quantum chaos and the absence of quasiparticles}\label{sec:chaos}

\subsection{Quantum chaos}

\subsubsection{Operator growth}

A quantum many-body system has no classical phase-space trajectory in the usual sense. To measure sensitivity to initially simple perturbations, one studies operator growth. Let $W(t)=e^{iHt}We^{-iHt}$. If $W$ and $V$ initially act on effectively independent degrees of freedom, then $[W(0),V]$ may be small or vanish. Under chaotic evolution, $W(t)$ becomes increasingly nonlocal in operator space, and the commutator grows.

For self-adjoint even operators $W$ and $V$, a nonnegative diagnostic is $C(t)=-\langle[W(t),V]^2\rangle_\beta$. A more analytically controlled object is a regularized OTOC such as
\begin{equation*}
F(t)=\Tr\left(
\rho^{1/4}W(t)\rho^{1/4}V
\rho^{1/4}W(t)\rho^{1/4}V
\right).
\end{equation*}

\subsubsection{Exponential growth}

In a large-$N$ scrambling window one finds
\begin{equation}
F(t)=F_{\mathrm{disc}}-\frac{A}{N}e^{\lambda_L t}+\cdots.
\label{eq:OTOCgrowth}
\end{equation}
In the leading strong-coupling, large-$N$ approximation,
\begin{equation}
\lambda_L=\frac{2\pi}{\beta}=2\pi T.
\label{eq:maxchaos}
\end{equation}
The correction becomes order one when
\begin{equation*}
\frac1N e^{\lambda_L t_*}\sim1,
\end{equation*}
so
\begin{equation*}
t_*\sim \frac{\beta}{2\pi}\log N.
\end{equation*}
This estimate assumes that the prefactor $A$ is of order one in the chosen scaling. If $\beta$ grows with $N$, its temperature dependence must also be retained when locating the scrambling time.

\subsubsection{The chaos bound}

Maldacena, Shenker, and Stanford~\cite{MSS} derived, under assumptions of analyticity, factorization, and thermal regularity, the bound
\begin{equation}
\lambda_L\le \frac{2\pi}{\beta}.
\label{eq:chaosbound}
\end{equation}
SYK saturates it in the strong-coupling large-$N$ limit.

The relevant mathematical statement is not that every finite-$N$ OTOC equals a pure exponential. The claim is an asymptotic one with a time window: $\beta\ll t\ll \beta\log N$. A rigorous maximal-chaos theorem would need uniform control of errors in precisely such a joint limit.

\subsubsection{Retarded ladder kernel}

The exponent $\lambda_L$ can be extracted from an analytically continued ladder equation. One looks for eigenfunctions of a retarded kernel with time dependence
\begin{equation*}
f(t_1,t_2)\sim e^{\frac{\lambda}{2}(t_1+t_2)}\phi(t_1-t_2).
\end{equation*}
The leading value of $\lambda$ for which the corresponding eigenvalue equals one gives the Lyapunov exponent. The same exceptional reparametrization sector that produces the $h=2$ Euclidean mode controls the maximal real-time growth.

Thus the same sector relates three calculations:
\begin{equation*}
\text{nearly zero Euclidean mode}
\quad\leftrightarrow\quad
\text{Schwarzian action}
\quad\leftrightarrow\quad
\text{maximal OTOC growth}.
\end{equation*}

\subsection{Fermi liquids}

The term ``non-Fermi liquid'' is meaningful only after the Fermi-liquid paradigm is stated carefully.

\subsubsection{Retarded Green function near a Fermi surface}

Consider an interacting translationally invariant fermion system in spatial dimension $d\ge2$. The retarded Green function is
\begin{equation}
G^R(\omega,k)
=\frac1{\omega-\varepsilon_{k}-\Sigma^R(\omega,k)}.
\label{eq:retarded}
\end{equation}
A Fermi surface may be characterized by a singularity or zero of the inverse Green function at zero frequency. Near a regular point $k_F$, let $k_\perp$ be the momentum normal to the Fermi surface.

Landau Fermi-liquid theory predicts
\begin{equation}
G^R(\omega,k)
\simeq
\frac{Z}{\omega-v_F^*k_\perp+i\Gamma(\omega,T)}
+G^R_{\mathrm{inc}},
\label{eq:FLpole}
\end{equation}
where $Z>0$ is the quasiparticle residue.

\subsubsection{The lifetime criterion}

A pole at $z_{k}=E_{k}-i\Gamma_{k}$ represents a useful quasiparticle only if its width is asymptotically small compared with its energy:
\begin{equation}
\frac{\Gamma_{k}}{|E_{k}|}\longrightarrow0
\quad\text{as }k\to k_F.
\label{eq:qpcriterion}
\end{equation}
For an ordinary Fermi liquid, $\Gamma(\omega,T)\sim c_1\omega^2+c_2T^2$ up to dimension-dependent logarithms. Hence $\Gamma(\omega,0)=o(|\omega|)$. This is the core mathematical content of the statement that low-energy quasiparticles are long lived.

The quasiparticle residue is formally
\begin{equation*}
Z=
\left[
1-\partial_\omega\operatorname{Re}\Sigma^R(\omega,k_F)|_{\omega=0}
\right]^{-1},
\end{equation*}
and also appears as the jump of the momentum distribution at the Fermi surface in a standard Fermi liquid.

\begin{definition}[Quasiparticle-like singularity]
A low-energy state is quasiparticle-like at $k_F$ if there exist $Z_*>0$, $v_F^*\ne0$, and $\Gamma(\omega)=o(|\omega|)$ such that the most singular part of $G^R$ has the form
\begin{equation*}
\frac{Z_*}{\omega-v_F^*k_\perp+i\Gamma(\omega)}.
\end{equation*}
\end{definition}

This definition captures the aspect of Fermi-liquid theory most relevant to SYK: the existence of an asymptotically sharp one-particle pole.

\subsection{SYK and non-Fermi-liquid behavior}

\subsubsection{A Green-function formulation}

The label ``non-Fermi liquid'' covers several mechanisms. A useful sufficient signal is that the self-energy dominates the bare frequency in the infrared. Suppose $\Sigma^R(\omega)\sim c(-i\omega)^\alpha, \qquad 0<\alpha<1$. Then
\begin{equation*}
\frac{|\Sigma(\omega)|}{|\omega|}
\sim |\omega|^{\alpha-1}\to\infty.
\end{equation*}
The interaction correction is larger than the bare kinetic term, and a narrow quasiparticle pole cannot be treated perturbatively.

\subsubsection{SYK scaling}

For quartic SYK,
\begin{equation*}
G(\tau)\sim \frac{\sgn\tau}{|\tau|^{1/2}}.
\end{equation*}
Fourier transformation gives, schematically, $G(i\omega) \sim i\,\sgn(\omega)|\omega|^{-1/2}$. Since in the infrared $G^{-1}\simeq-\Sigma$, we obtain $\Sigma(i\omega) \sim i\,\sgn(\omega)|\omega|^{1/2}$. Therefore
\begin{equation*}
\frac{|\Sigma(\omega)|}{|\omega|}
\sim |\omega|^{-1/2}\to\infty.
\end{equation*}
The self-energy therefore dominates the bare frequency in this approximation.

Another way to say the same thing is that the retarded Green function has a branch-type power-law singularity rather than a simple isolated quasiparticle pole: $G^R(\omega)\sim (\omega+i0^+)^{-1/2}$ up to a complex phase fixed by causality and particle-hole symmetry.

\subsubsection{The meaning of non-Fermi liquid}

The basic SYK model is zero-dimensional in space. It has no momentum and therefore no literal Fermi surface. It is more precise to call the infrared state a local non-Fermi-liquid fixed point or a non-quasiparticle quantum liquid. To describe an actual metal one must add charge and spatial structure.

\subsubsection{Complex SYK}

Introduce complex fermions $c_i,c_i^\dagger$ and a conserved charge
\begin{equation*}
Q=\sum_i c_i^\dagger c_i.
\end{equation*}
A representative Hamiltonian, with couplings antisymmetric in each pair and satisfying $g_{ij;kl}=\overline{g_{kl;ij}}$, is
\begin{equation*}
H=\sum_{i,j,k,l}g_{ij;kl}c_i^\dagger c_j^\dagger c_k c_l
-\mu\sum_i c_i^\dagger c_i.
\end{equation*}
In this paragraph only, use the conventional complex-fermion function $G(\tau)=-\langle T_\tau c_i(\tau)c_i^\dagger(0)\rangle$. With coupling variance scaled as $U^2/N^3$ and combinatorial constants absorbed into $U$, the large-$N$ equations take the form
\begin{equation*}
G(i\omega_n)=\frac1{i\omega_n+\mu-\Sigma(i\omega_n)},
\end{equation*}
 $\Sigma(\tau)=-U^2G(\tau)^2G(-\tau)$, together with a density constraint~\cite{NFLReview}. The preceding Majorana sections use the opposite Euclidean sign convention.

\subsubsection{SYK lattices and strange metals}

One can put SYK dots on a lattice, $x\in\Lambda$, with random interactions within each site and hopping between sites. A schematic Euclidean action is
\begin{align*}
L&=\sum_{x,i}c_{i,x}^\dagger(\partial_\tau-\mu)c_{i,x}
-\sum_{\langle xx'\rangle,ij}t_{ij,xx'}c_{i,x}^\dagger c_{j,x'}\\
&\quad
-\sum_{x,ijkl}g_{ijkl,x}
 c_{i,x}^\dagger c_{j,x}^\dagger c_{k,x}c_{l,x}.
\end{align*}
Such models can have incoherent transport without a quasiparticle Boltzmann equation.

A ``strange metal'' is a phenomenological term, often associated with a broad regime of approximately linear resistivity, $\rho(T)=\rho_0+AT$, and relaxation times of order the thermal time
\begin{equation*}
\tau_{\mathrm{rel}}\sim\frac{\hbar}{k_BT}.
\end{equation*}
These are related to, but not synonymous with, non-Fermi-liquid behavior.

\subsection{Transport and spatial extensions}

The physical importance of SYK is not that ordinary materials literally contain all-to-all Gaussian random couplings. Its importance is that it gives a controlled non-quasiparticle fixed point which can be embedded into more realistic constructions.

\subsubsection{A solvable strongly interacting fixed point}

Most analytic many-body calculations start from weakly interacting particles. SYK instead starts from a strongly interacting limit in which $|\Sigma(\omega)|\gg|\omega|$. The low-energy Green function and scaling dimension are still calculable. Thus it provides a rare explicit model of strong-coupling universality.

\subsubsection{Transport beyond kinetic theory}

Conventional transport is often derived from a Boltzmann equation $\partial_t f+v\cdot\nabla_x f=I_{\mathrm{coll}}[f]$. This presupposes a distribution of long-lived quasiparticles. In an SYK-derived incoherent metal, conductivity must instead be obtained directly from current-current correlation functions using a Kubo formula,
\begin{equation*}
\sigma(\omega)=\frac1{i\omega}G^R_{JJ}(\omega,k=0)
\end{equation*}
with the usual contact-term qualifications. This makes SYK a laboratory for transport when quasiparticle kinetic theory is unavailable.

\subsubsection{Hydrodynamics without quasiparticles}

Even without quasiparticles, conserved quantities can produce hydrodynamics. In spatial SYK models, energy and charge can diffuse: $\partial_t \varepsilon=D_E\nabla^2\varepsilon+\cdots$. The coexistence of local quantum criticality, diffusion, and chaos provides a controlled setting in which to study relations among diffusion constants, butterfly velocities, and thermal time scales.

\subsection{Random-matrix statistics}

SYK is a structured random matrix, not a Wigner matrix. The distinction matters.

\subsubsection{Global versus local spectral questions}

Let the eigenvalues be $E_1\le\cdots\le E_D$ with $D=2^{N/2}$. The empirical measure is
\begin{equation*}
\mu_N=\frac1D\sum_{a=1}^{D}\delta_{E_a/s_N}
\end{equation*}
for a suitable scale $s_N$. Global convergence of $\mu_N$ is a different question from local spacing statistics.

\subsubsection{Symmetry sectors and local statistics}
For quartic Majorana SYK, the antiunitary symmetries give the following random-matrix classes within a fixed fermion-parity sector~\cite{Cotler}:
\begin{center}
\begin{tabular}{@{}ccl@{}}
\toprule
$N\bmod8$ & class & symmetry qualification \\
\midrule
$0$ & GOE & antiunitary symmetry squares to $+1$ within a sector \\
$2,6$ & GUE & antiunitary symmetry exchanges the parity sectors \\
$4$ & GSE & antiunitary symmetry squares to $-1$ within a sector \\
\bottomrule
\end{tabular}
\end{center}
GOE, GUE, and GSE denote the Gaussian orthogonal, unitary, and symplectic ensembles. For $N\equiv4\pmod8$, count each Kramers pair once when computing spacings; for $N\equiv2,6\pmod8$, do not combine the paired parity spectra. Otherwise exact degeneracies create zero gaps unrelated to level repulsion. The table concerns $q=4$; other interaction orders require their own symmetry analysis.

Fix a sector and remove these multiplicities. If $E_*$ is a bulk energy with smooth mean density $\bar\rho_N(E_*)>0$, the locally unfolded energies are
\begin{equation}
x_a=\bar\rho_N(E_*)(E_a-E_*).
\label{eq:unfoldedLevels}
\end{equation}
This normalization gives mean spacing of order one when the density varies negligibly across the window. The random-matrix prediction concerns the point process $\sum_a\delta_{x_a}$, including its gap and correlation statistics. In the three Gaussian ensembles, the bulk spacing density behaves as $p(s)\propto s^r$ for small $s$, with $r=1,2,4$, respectively. Identifying a symmetry class does not prove convergence of the SYK point process.

At a spectral edge the mean density vanishes, so \eqref{eq:unfoldedLevels} is not the appropriate normalization. One instead centers the edge and rescales by a scale such as $\varepsilon_*$ in \eqref{eq:introEdgeScale}. A square-root density alone does not determine the edge point process. Bulk and edge universality are separate questions.

\subsubsection{Spectral form factor}

Define $K(\beta,t)=\E\left|Z(\beta+it)\right|^2$. In eigenvalue coordinates,
\begin{equation*}
K(\beta,t)=
\E\sum_{a,b}
 e^{-\beta(E_a+E_b)}e^{-it(E_a-E_b)}.
\end{equation*}
Thus $K$ is a Fourier transform of a two-level correlation measure. Chaotic random-matrix statistics produce the qualitative sequence
\begin{equation*}
\text{dip}\longrightarrow\text{ramp}\longrightarrow\text{plateau}.
\end{equation*}
The ramp is a long-range signature of spectral rigidity. Resolving finitely many local gaps alone does not establish a ramp: the latter also requires control of correlations over a growing range of energies and of the time scale in the Fourier transform.

\subsubsection{Different time scales}

In a semiclassical scrambling regime, $1\ll\beta J\ll N$, one distinguishes the thermal time $\beta$, the scrambling time, and the local Heisenberg time
\begin{equation*}
t_H(E)\asymp\delta_N(E)^{-1},
\qquad \delta_N(E):=\bar\rho_N(E)^{-1}.
\end{equation*}
The density here counts the levels of the chosen symmetry sector. The scrambling estimate $t_*\sim\beta(2\pi)^{-1}\log N$ further assumes an order-one OTOC prefactor, as discussed after \eqref{eq:OTOCgrowth}. The Heisenberg time depends on the energy window and is typically exponential in $N$ in a many-body spectral region. Neither time should be identified with the polynomial Schwarzian scale $C$.

\addtocontents{toc}{\protect\columnbreak}
\section[Nearly-AdS2 gravity]{The relation to nearly-$\AdS_2$ gravity}\label{sec:gravity}

\subsection{Geometric preliminaries}

We introduce the geometric objects needed to compare the SYK soft action with the boundary action of Jackiw--Teitelboim gravity. For fuller derivations, see~\cite{MSY,Trunin}.

\subsubsection{Lorentzian manifolds}

A spacetime is a smooth manifold $M$ equipped with a Lorentzian metric $g$ of signature $(-,+,\ldots,+)$. The metric gives the squared interval $ds^2=g_{\mu\nu}dx^\mu dx^\nu$. Vectors are timelike, null, or spacelike according as $g(v,v)$ is negative, zero, or positive. The Levi-Civita connection determines geodesics and the curvature tensor.

Einstein's equation is
\begin{equation}
R_{\mu\nu}-\frac12Rg_{\mu\nu}+\Lambda g_{\mu\nu}
=8\pi G_N T_{\mu\nu}.
\label{eq:Einstein}
\end{equation}
Anti-de Sitter space is a maximally symmetric vacuum solution with negative cosmological constant in dimension greater than two. The two-dimensional theory used below instead imposes curvature through a dilaton equation.

\subsubsection{$\AdS_2$}

Two-dimensional anti-de Sitter space can be written in Poincar\'e coordinates as
\begin{equation}
ds^2=\frac{L^2}{z^2}(-dt^2+dz^2),
\qquad z>0.
\label{eq:ads2}
\end{equation}
Its scalar curvature is constant,
\begin{equation*}
R=-\frac2{L^2}.
\end{equation*}
Its isometry group is locally $SL(2,\R)$. The same group is the global conformal group of one-dimensional time. This common symmetry is the first structural hint of the SYK-$\AdS_2$ connection.

\subsubsection{Horizons and black holes}

An event horizon is a global causal boundary: the boundary of the region from which future-directed causal curves fail to reach the asymptotic exterior. It is not a material membrane. In smooth large black holes, local curvature at the horizon may be small.

For the SYK story, charged near-extremal black holes are especially important. Extremality means that, at fixed conserved charge, the black hole reaches the minimum allowed mass and has zero Hawking temperature in the classical solution. Near extremality, the near-horizon geometry develops a long approximately $\AdS_2$ throat.

\subsection{JT gravity and the Schwarzian}

\subsubsection{Why two-dimensional gravity is special}

In two dimensions, the Einstein-Hilbert integral
\begin{equation*}
\int_M\sqrt g\,R\,d^2x
\end{equation*}
is topological up to boundary terms by Gauss-Bonnet. Pure two-dimensional Einstein gravity therefore has no local propagating metric degree of freedom.

Jackiw-Teitelboim gravity introduces a scalar dilaton $\Phi$. In Euclidean signature a common normalization is
\begin{align}
I_{JT}[g,\Phi]
&=-\frac{\Phi_0}{16\pi G}\int_M\sqrt g\,R\,d^2x\\
&\quad
-\frac1{16\pi G}\int_M\sqrt g\,\Phi\left(R+\frac2{L^2}\right)d^2x
+I_{\partial M}.
\label{eq:JT}
\end{align}
Varying with respect to $\Phi$ gives
\begin{equation*}
R+\frac2{L^2}=0.
\end{equation*}
Thus the metric is locally $\AdS_2$. The nontrivial degree of freedom is encoded in the dilaton and in the boundary embedding.

\subsubsection{Boundary reparametrization}

Regulate the asymptotic boundary by a curve parametrized by $u$. The physical boundary time is related to Poincar\'e time by $t=f(u)$. Solving the bulk constraints and evaluating the boundary action gives
\begin{equation}
I_{\mathrm{grav}}[f]
=-C\int_0^\beta du
\left\{
\tan\left(\frac{\pi f(u)}{\beta}\right),u
\right\}.
\label{eq:gravsch}
\end{equation}
This is the same functional form as the SYK soft action \eqref{eq:SchAction}.

The safest mathematical statement of the low-energy duality is therefore $I_{\mathrm{SYK,soft}}[f] \cong I_{\mathrm{JT,boundary}}[f]$ after matching the coefficient $C$ to $N\alpha_S$.

\subsubsection{Symmetry interpretation}

On the SYK side, the conformal equations possess approximate $\Diff(S^1)$ covariance, but a saddle preserves only $PSL(2,\R)$ and the UV kinetic term explicitly breaks the larger symmetry.

On the gravity side, large boundary reparametrizations are physical except for the exact $PSL(2,\R)$ isometries of $\AdS_2$. The Schwarzian measures the action cost of these boundary modes.

Thus the matching is not merely numerical. It is a matching of symmetry-breaking structures: $\Diff(S^1)/PSL(2,\R)$ appears on both sides.

\subsection{Black-hole thermodynamics}

\subsubsection{Bekenstein-Hawking entropy}

For a semiclassical black hole,
\begin{equation*}
S_{BH}=\frac{k_Bc^3}{4G_N\hbar}\mathcal A_H.
\end{equation*}
In units $c=\hbar=k_B=1$,
\begin{equation}
S_{BH}=\frac{\mathcal A_H}{4G_N}.
\label{eq:BHentropy}
\end{equation}
For a near-extremal charged black hole at fixed charge,
\begin{equation}
S_{BH}(T)=S_0+\gamma_{BH}T+O(T^2).
\label{eq:BHnear}
\end{equation}
Compare with large-$N$ SYK, $S_{SYK}(T)=Ns_0+N\gamma_{SYK}T+\cdots$. The functional form is identical.

This thermodynamic parallel predates the detailed Schwarzian matching and was an important clue that the same low-energy effective theory might be present.

\subsubsection{Residual entropy and order of limits}

On the black-hole side, $S_0$ is proportional to the extremal horizon area. On the SYK side, $Ns_0$ is the large-$N$ zero-temperature entropy obtained when $N\to\infty$ is taken before $T\to0$.

The correspondence should not be interpreted as an exponentially degenerate exact finite-$N$ SYK ground state. Instead, the density of low-energy states becomes exponentially large in a scaling window whose spacing is invisible before taking the thermodynamic limit.

\subsubsection{Schwarzian density of states}

After matching $C$, the JT boundary theory gives the same Laplace density \eqref{eq:SchDensity} and partition function \eqref{eq:SchPartition}. Restoring the edge energy and entropy gives \eqref{eq:SchMatching}. This compares smooth low-energy thermodynamics; microscopic level correlations require information beyond this density.

\subsection{The holographic dictionary}

It is helpful to phrase holography as an inverse problem for correlation data.

\subsubsection{Boundary dimensions and bulk masses}

A scalar field of mass $m$ in $\AdS_2$ is related to a boundary operator of dimension $h$ by
\begin{equation}
m^2L^2=h(h-1).
\label{eq:massdimension}
\end{equation}
SYK provides a tower of bilinear dimensions $h_n$ through $k(h_n)=1$. One may therefore assign a candidate bulk field $\phi_n$ with $m_n^2L^2=h_n(h_n-1)$. A schematic bulk action is
\begin{equation*}
\mathcal L_{\mathrm{bulk}}
=\sum_n\left[
\frac12(\partial\phi_n)^2+\frac12m_n^2\phi_n^2
\right]
+\frac1{\sqrt N}\sum_{nmk}\lambda_{nmk}\phi_n\phi_m\phi_k+\cdots.
\end{equation*}
Boundary three-point coefficients determine the cubic couplings after known $\AdS_2$ normalization factors are taken into account.

\subsubsection{What is established and what is not}

The robust statement is that the universal low-energy reparametrization sector of SYK is described by the same Schwarzian action as nearly-$\AdS_2$ JT gravity. This yields matching low-temperature thermodynamics and matching soft contributions to correlation functions and chaos.

A much stronger statement would be an exact microscopic equivalence $\cH_{\SYK}\cong\cH_{\mathrm{pure\ JT}}$ with equality of all finite-$N$ observables. That is not what has been established. Ordinary SYK also contains an infinite tower of order-one-dimension bilinear operators, so pure JT gravity captures only the universal soft sector, not all microscopic data.

\begin{remark}[A formulation of effective holographic matching]
For a specified class of low-energy sources $g_{\partial}$ and a prescribed joint limit $N\to\infty$, $\beta\to\infty$, the generating functional of the soft sector satisfies
\begin{equation*}
W^{\mathrm{soft}}_{\SYK}[g_{\partial}]
-W^{\mathrm{boundary}}_{JT}[g_{\partial}]
\to0
\end{equation*}
with controlled errors.
\end{remark}

To turn this template into a precise conjecture, one must specify the sources, normalization, joint limit, and topology of convergence.

\subsection{Black holes and chaos}

Near a black-hole horizon, a small perturbation sent in at an early time undergoes an exponentially large relative boost compared with a later probe. In the semiclassical shock-wave calculation this produces an exponential effect with scale
\begin{equation*}
\lambda_L=\frac{2\pi}{\beta}.
\end{equation*}
The same exponent appears from the Schwarzian mode and from the SYK retarded ladder kernel.

Thus the triangle
\begin{equation*}
\SYK\quad\longleftrightarrow\quad\text{Schwarzian}\quad\longleftrightarrow\quad\text{near-extremal black hole}
\end{equation*}
explains the common value of the chaos exponent. The common effective reparametrization theory accounts for this agreement.

Restoring constants,
\begin{equation*}
\lambda_L=\frac{2\pi k_BT}{\hbar},
\end{equation*}
so the Lyapunov time is
\begin{equation*}
\tau_L=\frac{\hbar}{2\pi k_BT}.
\end{equation*}
The same thermal time also appears in discussions of Planckian relaxation in strange metals, although the equality of all such time scales is not a universal theorem.

\subsection{Complex SYK and charge}

The charged complex SYK model has an additional parameter: the conserved density
\begin{equation*}
\mathcal Q=\frac1N\langle Q\rangle.
\end{equation*}
At nonzero density, the low-frequency spectral function need not be particle-hole symmetric. The conformal Green function contains a spectral-asymmetry parameter, often denoted $\mathcal E$ or an equivalent angle.

A notable thermodynamic relation links spectral asymmetry to the charge derivative of the zero-temperature entropy:
\begin{equation*}
\frac{\partial s_0}{\partial \mathcal Q}
\propto \mathcal E
\end{equation*}
with convention-dependent proportionality constants. An analogous relation appears in the thermodynamics of charged near-extremal black holes with an $\AdS_2$ throat.

The comparison therefore involves charge response and spectral asymmetry in addition to the entropy~\cite{SachdevBH,SachdevReview}.

\section[Probability and further questions]{Probabilistic results and further questions}\label{sec:probability}

\subsection{The probabilistic formulation}

We now recast SYK in a language closer to random operators and mean-field disordered systems.

\subsubsection{A sequence of random operator models}

For each even $N$, let $(\Omega_N,\mathscr F_N,\Pp_N)$ carry independent standard normal variables $g_I$. Let $\cH_N\cong\C^{2^{N/2}}$ and define
\begin{equation*}
H_N(\omega)=\sigma_N\sum_{I\in\mathcal I_N}g_I(\omega)\psi_I,
\end{equation*}
where $\psi_I=i^{q/2}\psi_{i_1}\cdots\psi_{i_q}$ and $\sigma_N^2=(q-1)!/N^{q-1}$.
The quenched pressure is
\begin{equation*}
p_N(\beta)=\frac1N\E\log\Tr e^{-\beta H_N}.
\end{equation*}
A basic mathematical problem is to prove existence and characterize $p(\beta)=\lim_{N\to\infty}p_N(\beta)$. Gamarnik, Pernice, Schmidhuber, and Zlokapa~\cite{Gamarnik} establish the annealed and quenched limits at sufficiently high fixed temperature. The September 2026 revision of He~\cite{He} reports convergence to the Schwinger--Dyson pressure for every fixed $\beta>0$ and fixed even $q\geq4$, as well as the limiting spectral edge. These thermodynamic statements do not themselves supply uniform infrared or real-time correlation estimates.

\subsubsection{Fluctuations at fixed temperature}
Write $P_N(\beta,g)=N^{-1}\log Z_g(\beta)$. Beyond its deterministic limit, one may ask for a normalization $a_N$ and a nondegenerate limit
\begin{equation*}
a_N\bigl(P_N(\beta,g)-\E P_N(\beta,g)\bigr)
\Longrightarrow X_\beta.
\end{equation*}
Gaussian concentration supplies an upper bound on variance, but neither determines $a_N$ nor establishes a Gaussian limit. The small-temperature-parameter expansion provides a useful check on possible scales.

\begin{example}[The first random Taylor coefficient]
Let $m_N=\binom Nq$ and $\tau_N=2^{-N/2}\Tr$. Since $\tau_N(H_q)=0$ and the Clifford monomials are orthogonal,
\begin{equation*}
P_N(0,g)=\tfrac12\log2,\qquad
\partial_\beta P_N(0,g)=0,\qquad
\partial_\beta^2P_N(0,g)=\frac{\sigma_N^2}{N}\sum_I g_I^2.
\end{equation*}
Consequently the centered second derivative $Y_N$ satisfies
\begin{equation*}
Y_N:=\partial_\beta^2P_N(0,g)-\E\partial_\beta^2P_N(0,g),
\qquad
\operatorname{Var}(Y_N)=\frac{2\sigma_N^4m_N}{N^2}
\sim\frac{2(q-1)!}{q}N^{-q}.
\end{equation*}
Moreover,
\begin{equation*}
\frac{N Y_N}{\sigma_N^2\sqrt{2m_N}}
=\frac1{\sqrt{2m_N}}\sum_I(g_I^2-1)
\Longrightarrow\mathcal N(0,1).
\end{equation*}
This is the ordinary central limit theorem for independent squared Gaussians. It concerns a Taylor coefficient at $\beta=0$, not the pressure at fixed positive $\beta$. Passing to the latter requires control of all remaining coefficients uniformly in $N$.
\end{example}

The annealed expansion
\begin{equation*}
\E Z_g(\beta)=\sum_{m=0}^\infty\frac{(-\beta)^m}{m!}\E\Tr H_q^m
\end{equation*}
has only even terms by symmetry of the disorder law. At fixed $N$, Gaussian integrability justifies the expansion. Its terms involve Gaussian pairings and Clifford signs. A fluctuation proof could use such expansions, sharper estimates on derivatives of $P_N$, or martingale differences obtained by revealing the couplings. In each approach, uniformity in $\beta$ is a separate issue.

\subsubsection{A dynamical order parameter}

The large-$N$ prediction is that
\begin{equation*}
G_\beta(\tau)
=\lim_{N\to\infty}
\frac1N\sum_i\E\langle T_\tau\psi_i(\tau)\psi_i(0)\rangle
\end{equation*}
is deterministic and solves
\begin{align*}
G^{-1}(i\omega_n)&=-\frac{i\omega_n}{2}-\Sigma(i\omega_n),\\
\Sigma(\tau)&=G(\tau)^3.
\end{align*}
This resembles a mean-field fixed-point equation, but the order parameter is a function rather than a scalar.

\subsubsection{Cavity intuition}

For the quartic model, separate one Majorana while keeping the $N$-dependent normalization fixed: $H_N=H_N^{(0)}+\psi_N B_N$, where
\begin{equation*}
B_N=\sigma_N\sum_{i<j<k<N}g_{ijkN}\psi_i\psi_j\psi_k,
\end{equation*}
and $H_N^{(0)}$ is the sum of the terms whose index sets do not contain $N$. Notice that $H_N^{(0)}$ uses $\sigma_N$, so it is not literally the normalized $(N-1)$-fermion Hamiltonian.
The distinguished mode $\psi_N$ interacts with a random dynamical bath generated by the remaining $N-1$ modes. In the large-$N$ limit, the bath's two-time correlations become self-consistent and deterministic. The resulting memory kernel is precisely the self-energy $\Sigma(\tau)=G(\tau)^3$. This is the noncommutative, dynamical analogue of a cavity field.

\subsubsection{Replica method}

The quenched average contains $\E\log Z$. Physicists often use
\begin{equation}
\E\log Z
=\lim_{n\to0}\frac{\E Z^n-1}{n}.
\label{eq:replica}
\end{equation}
For positive integers $n$, $Z^n$ introduces $n$ replicas. One computes $\E Z^n$, analyzes saddle points, and then analytically continues in $n$. This procedure is not automatically rigorous because the analytic continuation from integers to $n=0$ is not unique without additional control.

In ordinary SYK the standard saddle is replica diagonal in the usual parameter regime, unlike the low-temperature Parisi structure of the Sherrington-Kirkpatrick spin glass. The contrast is physically important: random couplings do not imply spin-glass order by themselves. Fermionic anticommutation changes the structure of frustration and low-energy states.

\subsection{Scope of the available results}

A useful map of the subject is the following.

\subsubsection{Finite-$N$ mathematics}

At finite $N$, the following objects are completely well-defined:
\begin{equation*}
H_N,\quad Z_N(\beta),\quad \rho_{\beta,N},\quad G_N(\tau),\quad
\operatorname{spec}(H_N),\quad \OTOC_N(t).
\end{equation*}
They are finite-dimensional random operator quantities.

\subsubsection{Rigorous large-$N$ results}

Relevant rigorous results include global spectral laws~\cite{Feng}, combinatorial melonic dominance~\cite{Bonzom}, and free-energy limits~\cite{Gamarnik,He}. Proposition~\ref{prop:concentration} supplies an elementary concentration bound in the present normalization. The exact hypotheses and normalizations vary among papers. These results should be viewed as mathematical foundations for parts of the standard physics picture, not as a single theorem proving every infrared claim.

\subsubsection{Controlled but formal physics}

The chain
\begin{equation*}
H_N
\xrightarrow{N\to\infty}
(G,\Sigma)
\xrightarrow{\omega\to0}
G_c
\xrightarrow{\text{soft fluctuations}}
I_{\Sch}
\xrightarrow{\text{real time}}
\lambda_L=2\pi/\beta
\end{equation*}
is supported by mutually consistent diagrammatic, saddle-point, numerical, and effective-field-theory calculations. But obtaining all arrows as one rigorous theorem with uniform error bounds in a joint low-temperature and large-$N$ limit remains a substantial challenge.

\subsubsection{The holographic layer}

The equality of the Schwarzian effective action on the SYK and JT sides is a concrete calculational correspondence. The interpretation as a full microscopic quantum-gravity dual is more ambitious and should be stated separately.

\subsection{Further questions}

The model generates a long list of precise questions.

\subsubsection{Derive the bilocal variational principle directly}

Beyond the pressure limit, derive the generating functional with time-dependent sources directly from the finite-$N$ Gibbs state. Specify a regularized collective action and prove convergence to the physically selected Schwinger--Dyson saddle, including its source derivatives. This requires control of saddle selection and correlation functions. An unrestricted infimum over $G$ and $\Sigma$ is not justified: the multiplier contour and the logarithmic determinant need not define a real minimization problem.

\subsubsection{Rigorous infrared scaling}

For $q=4$, take $\beta_N\to\infty$ with $\beta_N/N\to0$, and write $G_{N,\beta}$ for the averaged correlator in \eqref{eq:GN}, with its temperature dependence explicit. For every fixed $0<a<1/2$, a possible target is
\begin{equation}
\sup_{u\in[a,1-a]}
\left|\frac{G_{N,\beta_N}(\beta_Nu)}{G_{\beta_N,c}(\beta_Nu)}-1\right|
\longrightarrow0.
\label{eq:uniformConformalTarget}
\end{equation}
The relative normalization matters because both correlators tend to zero. An error estimate should distinguish the kinetic correction, which decreases at low frequency, from soft fluctuations governed by $\beta_N/C_N$. Their precise rates depend on the observable and norm; the scale comparison alone does not prove an $O(\beta_N^{-1})+O(\beta_N/N)$ bound. A theorem for times approaching the endpoints would require additional ultraviolet control.

\subsubsection{Derive the Schwarzian as a fluctuation theory}

At $\beta_N/C_N\to\lambda\in(0,\infty)$, formula~\eqref{eq:SchRescaled} suggests a normalized effective average with weight $e^{-\mathcal A[g]/\lambda}$ on the projective quotient. Specify its measure, a regularization, and observables, then prove convergence of the microscopic generating functional to those averages. One must control the remaining bilocal directions and their contribution to the coefficient. The Grassmann formulation does not supply an ordinary probability law on $G$ or $g$ without an additional construction.

At $\beta_N/C_N\to\infty$, the corresponding target is the quantum Schwarzian spectral representation, not a Gaussian expansion about $g(u)=u$. Correlation functions require matrix-element data as well as eigenvalue density~\cite{MertensTuriaciVerlinde}. For example, a self-adjoint finite-dimensional observable has the exact expansion
\begin{equation*}
\langle\mathcal O(\tau)\mathcal O(0)\rangle_{\beta,g}
=\frac1{Z_g(\beta)}\sum_{a,b}
e^{-(\beta-\tau)E_a-\tau E_b}|\mathcal O_{ab}|^2,
\qquad 0<\tau<\beta.
\end{equation*}
A Schwarzian limit must therefore identify the limit of the weighted two-energy measure appearing in this expression.

\subsubsection{Four-point convergence}

Prove convergence of connected four-point functions after $N$ rescaling:
\begin{equation*}
N\,\cF_N\longrightarrow (I-K)^{-1}\cF_{0,\infty}
\end{equation*}
in a suitable function space, where $\cF_{0,\infty}=\lim_{N\to\infty}N\cF_{0,N}$ is the rescaled seed. Control the singular $h=2$ sector by a separate perturbation theory.

\subsubsection{A maximal-chaos theorem}

Specify operators $V_N,W_N$, temperature scaling, and a time window such that
\begin{equation*}
C_N(t)=\frac{a}{N}e^{2\pi t/\beta}(1+o(1))
\end{equation*}
uniformly for $1\ll t/\beta\ll \log N$. This requires precise real-time analytic continuation and error control in a regime where small errors can be exponentially magnified.

\subsubsection{A mesoscopic spectral law}
Let $w$ be a nonnegative smooth function of compact support with $\int_\R w=1$. Smooth the unnormalized finite-$N$ counting measure by
\begin{equation*}
\rho_{N,\eta}(E):=\sum_{a=1}^{2^{N/2}}
\frac1\eta w\left(\frac{E-E_a}{\eta}\right).
\end{equation*}
This differs from the normalized global measure $\mu_N$. Let $C_N\asymp N$, and specify an edge centering $E_{0,N}$ and an entropy normalization $S_{0,N}$. A candidate local-law target is
\begin{equation}
\sup_{x\in[a,b]}
\left|\frac{\rho_{N,\eta_N}(E_{0,N}+x/C_N)}{e^{S_{0,N}}C_N}
-\frac1{2\pi^2}\sinh(2\pi\sqrt{2x})\right|
\xrightarrow{\Pp}0,
\quad 0<a<b<\infty.
\label{eq:mesoscopicTarget}
\end{equation}
On this interval the predicted mean spacing is of order $e^{-S_{0,N}}/C_N$, so a necessary scale separation is
\begin{equation*}
\frac{e^{-S_{0,N}}}{C_N}\ll\eta_N\ll\frac1{C_N}.
\end{equation*}
These inequalities express many levels per smoothing window and fine resolution on the Schwarzian energy scale. They are not sufficient by themselves to prove \eqref{eq:mesoscopicTarget}. The centering may require a sample-dependent correction, whose size must be controlled; a leading order-$N$ edge theorem does not determine it to accuracy $o(C_N^{-1})$.

As $x$ approaches zero, the local spacing changes and the edge scale \eqref{eq:introEdgeScale} replaces the preceding estimate. Smoothing is useful for a density statement; discreteness alone does not obstruct weak convergence of suitably normalized measures.

\subsubsection{Local spectral universality}
In a fixed symmetry sector, prove convergence of the bulk process defined by \eqref{eq:unfoldedLevels} to the corresponding Gaussian-ensemble process. Formulate the edge problem separately, with its own centering and scale. Gap distributions, ratios, and finite-order correlation functions are possible observables, but convergence of one is not automatically equivalent to convergence of all. Spectral form factors require additional estimates on the range of energies and times. These questions concern the correlations left undetermined by a smooth Schwarzian density.

\subsubsection{Universality in the disorder law}

Replace the standard normal variables $g_I$ by centered, unit-variance variables with sufficient moments. Ask whether the thermodynamic limit, infrared scaling, Schwarzian sector, and chaos exponent are universal. At the global spectral level this resembles Lindeberg replacement. At the level of soft modes and OTOCs it is far subtler.

\subsubsection{Sparse SYK}

Retain only a random subset of $q$-body interactions. The interaction structure becomes a random $q$-uniform hypergraph. Natural questions include the connectivity threshold for melonic mean-field behavior, concentration, localization, and the transition from sparse-cluster physics to dense SYK universality.

\subsubsection{Double scaling}

Let $q=q(N)$ grow with $N$. Then chord diagrams and noncommutative moment methods become natural. The fixed-$q$ and growing-$q$ limits give different asymptotic theories.

\subsubsection{Mathematical definition of a non-Fermi liquid}

For a spatial fermionic system, formulate hypotheses under which the failure of the quasiparticle property can be stated purely in terms of singularities of $G^R(\omega,k)$. A possible criterion is the absence of any decomposition
\begin{equation*}
G^R(\omega,k)
=\frac{Z_*}{\omega-v_F^*k_\perp+i\Gamma(\omega)}+R(\omega,k)
\end{equation*}
with $Z_*>0$, $\Gamma(\omega)=o(|\omega|)$, and $R$ less singular near the Fermi surface.

\subsubsection{Holography as an inverse problem}

Given conformal data $\{h_n,C_{nmk},C_{nmkl},\ldots\}$, find hypotheses ensuring the existence of a local or approximately local two-dimensional bulk action reproducing the boundary correlators order by order in $1/N$. Determine when such an action is unique and how bulk locality is encoded in asymptotics of boundary operator data.

\appendix
\section{Supplementary calculations}\label{app:identities}

\subsection{Collective-field calculations}

We record the coefficient and sign calculations used in the main text.

\subsubsection{From disorder to a quartic bilocal interaction}

For $q=4$, start from the Euclidean action
\begin{equation*}
S[\psi;g]
=\frac14\sum_i\int d\tau\,\psi_i\partial_\tau\psi_i
+\sigma_N\sum_{I\in\mathcal I_N}g_I\int d\tau\,\psi_I.
\end{equation*}
The Gaussian identity $\E e^{-gX}=e^{X^2/2}, \qquad g\sim\mathcal N(0,1)$, produces after disorder averaging a two-time interaction. Here $\sigma_N^2/2=3/N^3$ and the conversion of the sum over increasing quadruples to unrestricted flavor sums contributes $1/4!$. Thus, to leading order in $N$,
\begin{equation*}
\exp\left[
\frac N8\int d\tau_1d\tau_2
G(\tau_1,\tau_2)^4
\right]
\end{equation*}
with a positive sign in the Euclidean convention of these notes. The six exchanges needed to group the two sets of four Grassmann variables give sign $+1$.

Indeed, its coefficient is
\begin{equation*}
\frac{\sigma_N^2}{2}\frac{N^4}{4!}
=\frac{3}{N^3}\frac{N^4}{24}
=\frac N8.
\end{equation*}
This counting is the reason for the variance scaling $N^{-3}$ in quartic SYK and fixes the $G^4$ coefficient in \eqref{eq:collective}.

\subsubsection{From the collective action to Dyson}

Write $A=\tfrac12\partial_\tau-\Sigma$, and regularize the time variable so that traces are finite-dimensional. Since $\delta A=-\delta\Sigma$,
\begin{equation*}
\delta_\Sigma\left(-\frac12\log\det A\right)
=\frac12\Tr(A^{-1}\delta\Sigma)
=-\frac12\int (A^{-1})_{12}\,\delta\Sigma_{12}\,d\tau_1d\tau_2.
\end{equation*}
The last equality uses antisymmetry of the kernels. Adding the variation of $\tfrac12\int\Sigma G$ gives $G=A^{-1}$. Variation with respect to $G$ gives $\Sigma=G^{q-1}$. In particular, for $q=4$,
\begin{equation*}
G(i\omega_n)=\frac1{-i\omega_n/2-\Sigma(i\omega_n)},
\qquad \Sigma(\tau)=G(\tau)^3.
\end{equation*}
These are stationary equations; this calculation does not select the integration contour or prove dominance of a saddle.

\subsubsection{From scaling to $\Delta=1/q$}

Assume under $\tau\mapsto \lambda\tau$, $G(\lambda\tau)=\lambda^{-2\Delta}G(\tau)$. Then $\Sigma(\lambda\tau)=\lambda^{-2\Delta(q-1)}\Sigma(\tau)$. The convolution has scaling dimension $[\Sigma\circ G]=2\Delta q-1$, because one integrates once over time. To match $\delta(\tau)$, which has dimension $1$, require $2\Delta q-1=1$. Hence $\Delta=1/q$.

\subsubsection{From reparametrizations to the Schwarzian}

The strict infrared equations are invariant under \eqref{eq:reparamG}. Any reparametrized conformal solution $G_f$ has the same conformal action. The microscopic derivative term breaks this degeneracy. Locality in time, projective invariance, and the requirement of three derivatives at leading nontrivial order single out the Schwarzian derivative as the natural local functional. Thus one obtains
\begin{equation*}
I_{\mathrm{soft}}[f]
\propto N\int\{f,\tau\}\,d\tau
\end{equation*}
on the line, or the thermal-circle form \eqref{eq:SchAction}.

This argument explains the form of the action but not its coefficient. The coefficient $\alpha_S$ must be matched to the full theory, including the UV region neglected by the conformal approximation.

\subsection{The infrared Fourier transform}

For $0<\alpha<1$, the distributional Fourier transform obeys
\begin{equation*}
\int_0^\infty e^{i\omega t}t^{-\alpha}\,dt
=\Gamma(1-\alpha)e^{i\pi(1-\alpha)\sgn\omega/2}|\omega|^{\alpha-1}
\end{equation*}
where the integral is defined by inserting $e^{-\varepsilon t}$ and taking $\varepsilon\downarrow0$. For an odd function $f(t)=\sgn(t)|t|^{-2\Delta}$, using $\widehat f(\omega)=\int_\R e^{i\omega t}f(t)\,dt$, one obtains
\begin{equation*}
\widehat f(\omega)=2i\Gamma(1-2\Delta)\cos(\pi\Delta)
\sgn(\omega)|\omega|^{2\Delta-1},\qquad 0<2\Delta<1.
\end{equation*}
For $q=4$, $2\Delta=1/2$, hence $G(i\omega)\propto i\,\sgn(\omega)|\omega|^{-1/2}$. The time-domain power law thus gives the low-frequency singularity.

\subsection{The Clifford trace identity}

Let $\psi_I$ and $\psi_J$ be the Clifford monomials defined in \eqref{eq:Hqnormalized}. With normalized trace $\tau_N(A)=2^{-N/2}\Tr A$, one has orthogonality $\tau_N(\psi_I\psi_J)=0, \qquad I\ne J$, while $\tau_N(\psi_I^2)=1$. To verify orthogonality, let $X$ be a nonempty product of an even number of distinct generators and choose $\psi_k$ among its factors. Then $\psi_kX\psi_k=-X$, so cyclicity of trace gives $\Tr X=0$. For distinct $I,J$ of the same even size, $\psi_I\psi_J$ is a scalar multiple of such a product. The diagonal identity follows from $\psi_I^2=I$.

Consequently, for $H_q=\sum_IJ_I\psi_I$ with $J_I=\sigma_Ng_I$,
\begin{equation*}
\E\tau_N(H_q^2)
=\sum_I\E J_I^2
=\sigma_N^2\binom Nq.
\end{equation*}
In the present normalization this is exactly
\begin{equation*}
\E\tau_N(H_q^2)
=\frac Nq\prod_{r=1}^{q-1}\left(1-\frac rN\right)
\sim\frac Nq.
\end{equation*}
This gives a direct finite-dimensional explanation for the $O(N)$ second moment of the Hamiltonian.

\subsection[Schwarzian invariance]{Projective invariance of the Schwarzian}

Let
\begin{equation*}
M(f)=\frac{af+b}{cf+d},
\qquad ad-bc\ne0.
\end{equation*}
A direct differentiation gives
\begin{equation*}
(M\circ f)'=\frac{ad-bc}{(cf+d)^2}f'.
\end{equation*}
Substituting the first three derivatives into \eqref{eq:schwarzian}, all terms depending on $a,b,c,d$ cancel, yielding $\{M\circ f,\tau\}=\{f,\tau\}$. In particular,
\begin{equation*}
\left\{\frac{a\tau+b}{c\tau+d},\tau\right\}=0.
\end{equation*}
This identity is the elementary reason $PSL(2,\R)$ is the exact redundancy of the Schwarzian theory.

\subsection{Euclidean and real time}

For a finite-dimensional Hamiltonian with eigenvectors $\ket n$ and energies $E_n$, the thermal Euclidean correlator of fermionic operators has a Lehmann expansion. For $0<\tau<\beta$,
\begin{equation*}
G(\tau)
=\frac1Z\sum_{m,n}
e^{-\beta E_n}e^{-\tau(E_m-E_n)}
\langle n|\chi|m\rangle\langle m|\chi|n\rangle.
\end{equation*}
The retarded function is $G^R(t)=-i\theta(t)\langle\{\chi(t),\chi(0)\}\rangle_\beta$. Both are determined by the same spectral measure, and analytic continuation connects their frequency-domain representations. The difficulty in chaos calculations is not the finite-$N$ existence of analytic continuation, but obtaining uniform asymptotics after taking large $N$, low temperature, and long time.

\subsection{Comparison of conventions}

Different papers normalize Majorana fermions either by $\{\psi_i,\psi_j\}=\delta_{ij}$ or by $\{\psi_i,\psi_j\}=2\delta_{ij}$. The convention in these notes is the second one, so $\psi_i^2=I$. If $\chi_i:=\psi_i/\sqrt2$, then $\{\chi_i,\chi_j\}=\delta_{ij}I$ and
\begin{equation*}
i^{q/2}\psi_{i_1}\cdots\psi_{i_q}
=2^{q/2}i^{q/2}\chi_{i_1}\cdots\chi_{i_q}.
\end{equation*}
Thus the same operator $H_q$ has coupling variance $2^q(q-1)!/N^{q-1}$ when written in the $\chi$ convention. Correspondingly $G_\psi=2G_\chi$, while the equations used in these notes are
\begin{equation*}
G_\psi^{-1}(i\omega)=-\frac{i\omega}{2}-\Sigma_\psi(i\omega),
\qquad
\Sigma_\psi(\tau)=G_\psi(\tau)^{q-1}.
\end{equation*}
Likewise some authors write $G(\tau)=-\langle T_\tau\chi(\tau)\chi(0)\rangle$ with an explicit minus sign, while others absorb it into conventions. The universal structural statements are unaffected:
\begin{equation*}
\Sigma\propto J^2G^{q-1},
\qquad
\Delta=1/q,
\qquad
I_{\mathrm{soft}}\propto \frac NJ I_{\Sch},
\qquad
\lambda_L\to 2\pi/\beta.
\end{equation*}
When comparing formulas between papers, one should first align these conventions.

\clearpage
\section{Terminology}

\begin{longtable}{@{}>{\raggedright\arraybackslash}p{0.26\textwidth}>{\raggedright\arraybackslash}p{0.70\textwidth}@{}}
\toprule
Physics term & Mathematical translation \\
\midrule
\endfirsthead
\toprule
Physics term & Mathematical translation \\
\midrule
\endhead
Majorana fermion & Self-adjoint Clifford generator $\psi_i$ with $\{\psi_i,\psi_j\}=2\delta_{ij}I$. \\
Quenched disorder & Random Hamiltonian sampled once, then thermodynamics computed at fixed sample before averaging $\E\log Z$. \\
Self-energy & Kernel $\Sigma$ entering the nonlinear resolvent identity $G^{-1}=G_0^{-1}-\Sigma$. \\
Large $N$ & Asymptotic regime in which singlet fluctuations are suppressed and melonic diagrams dominate. \\
Melon & Recursively generated leading graph maximizing powers of the flavor index $N$. \\
Infrared & Frequency regime $|\omega|\ll1$, or time regime $|\tau|\gg1$, with the effective large-$N$ coupling scale set to one. \\
Emergent conformal symmetry & Approximate scale/reparametrization covariance of the low-energy Schwinger-Dyson equations after dropping the kinetic term. \\
Soft mode & Nearly zero-cost reparametrization fluctuation $f\in\Diff/PSL(2,\R)$. \\
Schwarzian & Effective action for the soft mode, built from $\{f,\tau\}$. \\
Quasiparticle & Pole-like excitation with width $\Gamma=o(|E|)$. \\
Non-Fermi liquid & Gapless fermionic state lacking the Landau quasiparticle pole structure. \\
OTOC & Four-point function probing growth of noncommutativity and operator complexity. \\
Maximal chaos & Exponential OTOC growth with $\lambda_L=2\pi/\beta$, saturating the chaos bound. \\
$\AdS_2$ & Constant negative curvature two-dimensional spacetime with isometry group $SL(2,\R)$. \\
JT gravity & Two-dimensional dilaton gravity whose boundary mode is governed by a Schwarzian action. \\
Holography & Reconstruction of bulk gravitational observables from lower-dimensional quantum correlation data in an appropriate limit. \\
\bottomrule
\end{longtable}

\clearpage
\phantomsection
\addcontentsline{toc}{section}{References}

\begin{thebibliography}{99}
\small
\setlength{\itemsep}{2pt}

\bibitem{Rosenhaus}
V. Rosenhaus,
\emph{An introduction to the SYK model},
\href{https://arxiv.org/abs/1807.03334}{arXiv:1807.03334} (2018).

\bibitem{SachdevYe}
S. Sachdev and J. Ye,
\emph{Gapless spin-fluid ground state in a random quantum Heisenberg magnet},
Phys. Rev. Lett. 70, 3339 (1993), \href{https://arxiv.org/abs/cond-mat/9212030}{arXiv:cond-mat/9212030}.

\bibitem{MaldacenaStanford}
J. Maldacena and D. Stanford,
\emph{Remarks on the Sachdev-Ye-Kitaev model},
Phys. Rev. D 94, 106002 (2016), \href{https://arxiv.org/abs/1604.07818}{arXiv:1604.07818}.

\bibitem{PolchinskiRosenhaus}
J. Polchinski and V. Rosenhaus,
\emph{The spectrum in the Sachdev-Ye-Kitaev model},
JHEP 04 (2016) 001, \href{https://arxiv.org/abs/1601.06768}{arXiv:1601.06768}.

\bibitem{KitaevSuh}
A. Kitaev and S. J. Suh,
\emph{The soft mode in the Sachdev-Ye-Kitaev model and its gravity dual},
JHEP 05 (2018) 183, \href{https://arxiv.org/abs/1711.08467}{arXiv:1711.08467}.

\bibitem{MSY}
J. Maldacena, D. Stanford, and Z. Yang,
\emph{Conformal symmetry and its breaking in two dimensional Nearly Anti-de-Sitter space},
Prog. Theor. Exp. Phys. 2016, 12C104, \href{https://arxiv.org/abs/1606.01857}{arXiv:1606.01857}.

\bibitem{MSS}
J. Maldacena, S. H. Shenker, and D. Stanford,
\emph{A bound on chaos},
JHEP 08 (2016) 106, \href{https://arxiv.org/abs/1503.01409}{arXiv:1503.01409}.

\bibitem{SachdevBH}
S. Sachdev,
\emph{Bekenstein-Hawking entropy and strange metals},
Phys. Rev. X 5, 041025 (2015), \href{https://arxiv.org/abs/1506.05111}{arXiv:1506.05111}.

\bibitem{SachdevReview}
S. Sachdev,
\emph{Quantum statistical mechanics of the Sachdev-Ye-Kitaev model and charged black holes},
\href{https://arxiv.org/abs/2304.13744}{arXiv:2304.13744}.

\bibitem{NFLReview}
D. Chowdhury, A. Georges, O. Parcollet, and S. Sachdev,
\emph{Sachdev-Ye-Kitaev models and beyond: a window into non-Fermi liquids},
Rev. Mod. Phys. 94, 035004 (2022), \href{https://arxiv.org/abs/2109.05037}{arXiv:2109.05037}.

\bibitem{GuLucasQi}
Y. Gu, X.-L. Qi, and D. Stanford,
\emph{Local criticality, diffusion and chaos in generalized Sachdev-Ye-Kitaev models},
JHEP 05 (2017) 125, \href{https://arxiv.org/abs/1609.07832}{arXiv:1609.07832}.

\bibitem{PatelSachdev}
A. A. Patel and S. Sachdev,
\emph{Theory of a Planckian metal},
Phys. Rev. Lett. 123, 066601 (2019), \href{https://arxiv.org/abs/1906.03265}{arXiv:1906.03265}.

\bibitem{Cotler}
J. S. Cotler et al.,
\emph{Black holes and random matrices},
JHEP 05 (2017) 118, \href{https://arxiv.org/abs/1611.04650}{arXiv:1611.04650}.

\bibitem{SSS}
P. Saad, S. H. Shenker, and D. Stanford,
\emph{JT gravity as a matrix integral},
\href{https://arxiv.org/abs/1903.11115}{arXiv:1903.11115}.

\bibitem{Bonzom}
V. Bonzom, V. Nador, and A. Tanasa,
\emph{Diagrammatic proof of the large $N$ melonic dominance in the SYK model},
Lett. Math. Phys. 109, 2611 (2019), \href{https://arxiv.org/abs/1808.10314}{arXiv:1808.10314}.

\bibitem{Feng}
R. Feng, G. Tian, and D. Wei,
\emph{Spectrum of SYK model},
\href{https://arxiv.org/abs/1801.10073}{arXiv:1801.10073}.

\bibitem{Anschuetz}
E. R. Anschuetz, C.-F. Chen, B. T. Kiani, and R. King,
\emph{Strongly interacting fermions are non-trivial yet non-glassy},
\href{https://arxiv.org/abs/2408.15699}{arXiv:2408.15699}.

\bibitem{Gamarnik}
D. Gamarnik, F. Pernice, A. Schmidhuber, and A. Zlokapa,
\emph{The free energy limit of the SYK model at high temperature},
\href{https://arxiv.org/abs/2605.02768}{arXiv:2605.02768}.

\bibitem{He}
Y. He,
\emph{Free energy and spectral edge of the SYK model},
\href{https://arxiv.org/abs/2607.18998v3}{arXiv:2607.18998v3} (2026).

\bibitem{StanfordWitten}
D. Stanford and E. Witten,
\emph{Fermionic localization of the Schwarzian theory},
JHEP 10 (2017) 008, \href{https://arxiv.org/abs/1703.04612}{arXiv:1703.04612}.

\bibitem{MertensTuriaciVerlinde}
T. G. Mertens, G. J. Turiaci, and H. L. Verlinde,
\emph{Solving the Schwarzian via the conformal bootstrap},
JHEP 08 (2017) 136, \href{https://arxiv.org/abs/1705.08408}{arXiv:1705.08408}.

\bibitem{Trunin}
D. A. Trunin,
\emph{Pedagogical introduction to SYK model and 2D dilaton gravity},
\href{https://arxiv.org/abs/2002.12187}{arXiv:2002.12187}.

\end{thebibliography}

\end{document}
